\documentclass[11pt,reqno]{amsart}
\usepackage[T1]{fontenc}
\usepackage{lmodern,amsmath,amssymb,mathtools,mathrsfs,microtype}
\usepackage[margin=1in]{geometry}
\usepackage{enumitem,booktabs,longtable,array}
\usepackage[hidelinks]{hyperref}
\hypersetup{pdftitle={General Theta Foundations I: Acquired Geometry and Causal Resource Transfer},pdfauthor={Qian Qi}}
\ifdefined\pdfinfoomitdate\pdfinfoomitdate=1\fi
\ifdefined\pdftrailerid\pdftrailerid{}\fi
\ifdefined\pdfsuppressptexinfo\pdfsuppressptexinfo=15\fi
\numberwithin{equation}{section}
\newtheorem{theorem}{Theorem}[section]
\newtheorem{proposition}[theorem]{Proposition}
\newtheorem{lemma}[theorem]{Lemma}
\newtheorem{corollary}[theorem]{Corollary}
\theoremstyle{definition}
\newtheorem{definition}[theorem]{Definition}
\newtheorem{example}[theorem]{Example}
\theoremstyle{remark}
\newtheorem{remark}[theorem]{Remark}
\newcommand{\E}{\mathbb E}
\newcommand{\Pp}{\mathbb P}
\newcommand{\R}{\mathbb R}
\newcommand{\N}{\mathbb N}
\newcommand{\one}{\mathbf 1}
\newcommand{\tr}{\operatorname{tr}}
\newcommand{\conv}{\operatorname{conv}}
\newcommand{\supp}{\operatorname{supp}}
\newcommand{\dist}{\operatorname{dist}}
\newcommand{\TV}{\mathrm{TV}}
\newcommand{\HS}{\mathrm{HS}}
\newcommand{\aut}{\mathrm{aut}}
\newcommand{\cp}{\mathrm{cp}}
\newcommand{\calR}{\mathcal R}
\newcommand{\calQ}{\mathcal Q}
\newcommand{\calC}{\mathcal C}
\newcommand{\calA}{\mathcal A}
\newcommand{\calF}{\mathcal F}
\newcommand{\ip}[2]{\langle #1,#2\rangle}
\newcommand{\norm}[1]{\lVert #1\rVert}
\newcommand{\law}{\operatorname{Law}}
\newcommand{\smin}{\mathop{\mathrm{ess\,inf}}}
\newcommand{\eps}{\varepsilon}
\newcommand{\bnorm}[1]{\lVert #1\rVert_{\mathrm b}}
\setlist{nosep,leftmargin=*}
\emergencystretch=2em
\title[Acquired geometry and causal resource transfer]{General Theta Foundations I:\\Acquired Geometry and Causal Resource Transfer}
\author{Qian Qi}
\date{October 2026}
\subjclass[2020]{Primary 60G40, 93E11; Secondary 28A80, 62C20, 81P15, 94A29}
\keywords{causal experiment, predictive quotient, finite-state observer, stopping time, occupation measure, acquired geometry}
\begin{document}
\begin{abstract}
We study causal prediction under simultaneously bounded expected acquisition and total finite-state memory. A prepared positive experiment with finite affine test closure gives an exact common-program occupation criterion and an actual-law quantization profile. We prove a budget-normalized perturbation theorem for the resulting constrained optimal risk: preparation and one-call errors enter through $\delta_0+N\eta$, without an additional horizon factor and without uniform absorption of the competing programs. The proof combines an exact survival envelope with a positive-map occupation estimate, and transfers each program at the same mean and state budgets. Its square-root order is optimal on the declared class. A singular preparation experiment yields an exact all-budget memory frontier and a prediction-risk law whose critical small-mass exponent is $(M-1)/M$, although the physical carrier has fixed dimension. Every lower bound ranges over all legal randomized machines. Classical Gaussian and noncommuting quantum instruments verify the same general transfer hypotheses; rank changes and rare acquired components retain their actual mass. Complete proofs of the underlying occupation, geometric, and finite-description constructions are included. Calibration, numerical precision, causal deficiency, simulator state, program descriptions, workspace, and physical time remain separately accounted for.
\end{abstract}
\maketitle
\section{Introduction}
The geometry relevant to a prediction problem is the geometry that an observer actually acquires. For a causal experiment this geometry depends on preparation, intervention and stopping, as well as on the physical kernel. A cover of a formally reachable set is not an acquisition procedure. Conversely, a procedure that realizes a favorable conditional distribution may consume many unsuccessful observations before producing its report. The central question is to determine the predictive information that can be acquired and retained by an executable observer, with these costs included.

We take the experiment, not a manifold or an information matrix, as primitive. A preparation and positive report instruments determine future-test evaluations and hence a predictive quotient. A finite observer has a fixed program: it selects an action from its current live label, reads the next report, and either changes live label or enters a terminal label. The latter decision is part of the program. The audit is performed only afterwards. The program cannot read the true model, its own unrecorded past, or an external clock. Persistent labels, program constants, temporary arithmetic, calibration and physical calls are different resources.

Three difficulties arise under an acquisition constraint $\E\tau\le N$. An algebraic balance equation can contain occupation never generated by the preparation. Optimizing the stopping rule changes the acquired law on which compression acts. Finally, approximation of a stopped law in a bounded task metric does not control its mean acquisition cost: a rare discrepancy can initiate an arbitrarily long excursion. A fixed-time quantization bound, or existence of an optimal Borel routing kernel, therefore does not by itself yield an executable finite-description observer under the original budget.

Theorem~\ref{thm:main} connects these three obligations. Its occupation and geometric parts are the causal-stopping results of~\cite{GTFMean}, included with their complete proofs. The constructive part of~\cite{GTFExactBudget}, reproduced with its proof, establishes that the exact optimum is the limit of certified finite programs using the \emph{same} mean and state budgets, without assuming that all competing observers are uniformly absorbing. This completion starts from uniformly effective unnormalized report instruments, not from an already computed optimum, a smooth routing rule, or a contraction of the normalized posterior. It applies to a known family of experiments with an unknown fixed true model; the program, including its readout, is common to the whole family.

The mechanism is a quantitative conversion of stopping freedom into budget slack. A small fresh probability of routing to an existing terminal label creates a strictly positive mean margin and changes bounded risk by at most that probability times the paid live occupation. No label, clock, or physical reset is added. An occupation-weighted resolvent estimate then separates terminal-risk error from expected-cost error. The former does not acquire an inverse stopping factor, whereas the latter does. This difference permits finite report averaging, probability rounding and certified evaluation to spend only part of the margin. Consequently the finite feasibility test may reject uncertain boundary cases and still contains an approximation to every admissible Borel competitor. The lower certificate is global, not an empirical failure of a search routine.

The new issue is the sharp dependence on the paid mean budget. A linear estimate of the stopping margin loses a factor of order $\sqrt N$. We instead optimize the entire survival sequence under its mean constraint. The resulting concave envelope gives a budget-normalized transfer: with $\zeta=\delta_0+N\eta$ and $c_N=1-N^{-1}$,
\[
 |\calR_1^\Theta(N,M)-\calR_2^\Theta(N,M)|
 \le L\{2\zeta+\sqrt{\zeta^2+2c_N\zeta}\}+\alpha,
\]
with truncation by $L$ and a separate linear one-call endpoint. Here $\delta_0$ and $\eta$ are preparation and raw one-call base-norm errors, and $\alpha$ is a common-task error. Both transfer directions preserve $N,M$. Thus a one-call accuracy of order $a^2/N$ suffices for risk accuracy of order $a$ when $N\ge2$. A bound on terminal-law error alone could not prove feasibility at the same budget.

The square-root order is genuinely necessary. A small preparation mass can activate a trapped component whose calls were absent under the nominal preparation. In a three-state raw experiment we compute the failure-risk frontier exactly for every $N\ge1$ and $M\ge2$, over all randomized autonomous machines. Its squared-prediction risk lies between one quarter and one half of that exact frontier. At a critical mean budget, for fixed $M$ and rare preparation mass $b\downarrow0$, the scale is
\[
 b^{(M-1)/M}.
\]
This is a joint acquisition--memory consequence of the actual trapped occupation, not a dimension assigned to the three physical coordinates. The $M=2$ frontier supplies a sharp preparation-perturbation witness; a second raw-kernel perturbation supplies the matching one-call-error witness. Their programs also work for an unknown fixed parameter in a prescribed interval. The lower bound uses a standard phase-type comparison in its precisely paid discrete-clock form, whose proof is included in Appendix~\ref{app:clock-order}.

For completeness, the inherited occupation argument uses one joint routing inequality over all models. Positivity removes every finite invariant occupation component without changing the terminal ensemble. The mean budget provides compactness, even when other programs never terminate. This proves attainment and finite-subfamily risk obstructions. These statements do not use a model-dependent program, an unrecorded clock, or a formal reachable distribution.

A sharp geometric consequence requires additional information about what a report reveals. Suppose the current task response is recoverable from the current report and the conditional one-call laws are comparable to a reference law $\nu$ that is itself physically prepared. Stopping with mean budget $N$ can increase the density of an acquired terminal law by at most a factor proportional to $N$. The optimal distortion under this density constraint is a mass-constrained quantization profile $\calQ_m^\nu(N)$. A rejection rule realizes the matching upper with one live state and $m$ terminal states. The independent-preparation frontier of the retained adapted-stopping article~\cite{GTFStopped}, reproduced here as a limiting case, gives the exact identities
\begin{equation}\label{eq:intro-exact}
 \calR_{\cp}(N,m)-B=\calQ_m^\nu(N),\qquad
 \calR_{\aut}(N,M)-B=\calQ_{M-1}^\nu(N).
\end{equation}
For correlated controlled reports the profile remains a two-sided comparison. If the acquired reference support is $s$-regular, it yields $(NM)^{-2/s}$ rather than a rate determined by the ambient parameter dimension. There is no requirement that $s$ be an integer. Atoms and rank changes are governed by the profile, not by an artificially assigned common dimension.

The distinction between the two levels is essential. The occupation theorem applies to informative normalized filters with a finite affine closure; the sharp profile theorem does not silently assume that such a filter is report-recoverable. Its classical realization is a genuinely correlated controlled observation experiment, and its quantum realization is a noncommuting instrument with continuously many outcomes. Neither realization is used in the proof of the general theorem. Known-kernel filtering, stationary observation, scheduled horizons and free stopping are separate tasks; a rate in one is not asserted in another without a causal comparison.

\subsection*{Relation to prior work}
Predictive states and minimal future-test representations have an established history; see Shalizi and Crutchfield~\cite{SC} and the experiment-comparison treatment of Fritz, Gonda, Perrone and Rischel~\cite{FGPR}. Finite belief compression is also classical, including the value-directed approach of Poupart and Boutilier~\cite{PB}. These works motivate, but do not supply, the hard-state mean-acquisition theorem proved here.

Occupation equations and constrained control are classical tools. In particular, Dufour and Prieto-Rumeau~\cite{DPR} characterize the invariant part of phantom occupation solutions. We do not claim that decomposition as new. Here it is used for a different purpose: a finite positive predictive representation, a report-only homogeneous controller, and terminal risk under an upper bound on all acquired occupation. One may retain a compact algebraic relaxation and remove its invariant part rather than assume uniform absorption. The proof must still enforce one routing program, including its stopping decision, across all models. Dufour and Prieto-Rumeau's later finite-mixture result~\cite{DPRG} uses uniform absorption for its occupation-space conclusions. Here the finite construction regularizes each competitor within its existing label set, proves a quantitative error against its actual occupation, and pays for the regularization inside the given mean budget. This is not a claim that occupation methods or termination regularization themselves are new. A recent preprint of Sokolov~\cite{Sokolov} gives a Markovization of stopping intensities while preserving observed and surviving measures. A policy of the form $\lambda(t,X_t)$ has access to a state and time; it is not, without further work, a controller with a prescribed finite number of persistent labels.

The functional $\calQ_m^\nu$ is normalized trimmed quantization, not a new static quantization objective. Cuesta-Albertos, Gordaliza and Matr\'an~\cite{CGM} established trimmed $k$-means, and V\'erine, Pydi, Negrevergne and Chevaleyre~\cite{VPNC} study budgeted rejection sampling for distributional objectives. The exact independent-preparation value and its one-live-state realization are inherited from~\cite{GTFStopped}, whose theorem already allows a nonconstant conditional Bayes risk and singular preparation. Those joint-routing and correlated common-program results are inherited from~\cite{GTFMean}. The exact-budget construction of~\cite{GTFExactBudget} turns the resulting Borel optimum into arbitrarily accurate finite-description certificates at exactly the same constraints; the present transfer sharpens its resource dependence and supplies converse witnesses. The small-ball and covering proof is included for every real acquisition budget, so no asymptotic quantization theorem is being invoked outside its hypotheses.

The retained adapted-stopping article~\cite{GTFStopped} also proves a distinct theorem for informative Gaussian and noncommuting weak-effect instruments: at each fixed positive acquisition price it matches the optimal full-history sensing-and-stopping value with a finite-state value. That proof uses calibrated occupation and cost absorption, not report recoverability. We retain it intact rather than identify a Lagrange price with a constrained mean budget without a duality theorem. Its error-composition and calibration mechanisms are inherited tools. Appendix~\ref{app:bayesian} likewise reproduces the earlier Bayesian information decomposition~\cite{GTFBayes}. The retained report-recoverable realizations supply explicit constrained curves. The informative Gaussian and quantum realizations of~\cite{GTFExactBudget}, retained with their proofs, verify the effective construction on overlapping posteriors, including nonmixing and changing-rank kernels. Their raw norm estimates also verify the new budget-normalized transfer, as stated in Corollary~\ref{cor:raw-normalized}. They provide certified optimal values and stability bounds, not an automatic closed-form power law for these broader classes.

Piunovskiy and Zhang~\cite{PZ} identify uniform absorption as the relevant condition for continuity of the strategic-to-occupation projection in their semicontinuous setting. Our conclusion concerns a different object: reoptimized bounded terminal risk after a same-budget change of the program. We do not claim continuity of unregularized occupation measures. Similarly, phase-type variability is classical: Aldous and Shepp~\cite{AS} prove the least-variance property, and O'Cinneide~\cite{OC} proves the stronger majorization comparison. The new singular-acquisition law comes from coupling that finite-phase limitation to the exact \emph{prepared} occupation and risk equations, not from claiming Erlang extremality anew.

Bayes-risk quantization and Bregman geometry are prior constructions~\cite{Varshney}. Kim's finite-memory study~\cite{Kim} compares fixed-policy costs under one closed-loop execution. Here the converse ranges over all declared adaptive acquisition and stopping policies, while the implementing rule is one common finite-state program. These different quantifiers, and not the mere use of a belief approximation, delimit the contribution.

\subsection*{Scope}
The paper proves a class theorem, not a finite-dimensional reduction of every causal experiment. Finite affine closure, common report domination and the audit interface are explicit in Section~\ref{sec:experiments}. Report recoverability and uniform comparison are additional hypotheses only for the sharp profile. There is no assumption of a preassigned Fisher matrix, no use of total variation to infer a large-deviation principle, and no promotion of finite matrices to unbounded-operator identities. An actual mean budget is not an almost-sure deadline. Under the effective raw certificates the exact constrained optimum has finite two-sided certificates. The sharp singular law is an additional consequence, not such a general informative-filter power law. Efficient synthesis, optimal program length, and a general intrinsic rate for informative filters remain further mathematical questions. The small-mass exponent is sharp for the stated machine class and protocol, not a claim about a free externally clocked or history-replaying observer.

\section{Prepared experiments and finite observers}\label{sec:experiments}
\subsection{The raw experiment and its predictive representation}
Let the hidden state, report and history spaces be standard Borel. A preparation and a finite action set $\calA$ specify normalized positive causal kernels
\[
 K_a^\theta(dx',dy\mid x),\qquad \theta\in\Theta.
\]
The model $\theta$ is fixed during a run and is not an input to the observer. The singleton case is included. Every chosen action incurs one acquisition. The observer may stop after any acquired report, but not before the first. After stopping, a specified protected audit is executed: it cannot be used to decide whether to continue. Fresh randomization is independent of the hidden experiment conditional on the visible history.

We impose the following finite predictive representation, to be verified from the raw kernels. Let $E$ be a finite-dimensional real vector space, $C\subset E$ a closed, pointed, generating cone, and $e\in E^*$ strictly positive on $C\setminus\{0\}$. Its base
\[
 D=\{x\in C:e(x)=1\}
\]
is compact. A preparation is $x_0^\theta\in D$. For each action, $Y_a$ is standard Borel and $\lambda_a$ is a finite positive measure. There are measurable positive linear maps
\begin{equation}\label{eq:raw-instrument}
 L_a^\theta(y):C\longrightarrow C,\qquad
 \int\norm{L_a^\theta(y)}\lambda_a(dy)<\infty,
 \qquad \int eL_a^\theta(y)x\,\lambda_a(dy)=e(x).
\end{equation}
The measures $\lambda_a$ are common to all models, but the integrability constants need not be uniform in $\theta$. Starting from $x$, the report has density $eL_a^\theta(y)x$ and the updated predictive state is
\begin{equation}\label{eq:update}
 F_a^\theta(y,x)=\frac{L_a^\theta(y)x}{eL_a^\theta(y)x}
\end{equation}
on positive-evidence reports. If the denominator is zero, positivity implies $L_a^\theta(y)x=0$; no positive-probability posterior is assigned there.

This representation means that the selected executable future tests, their report pullbacks and the audit moments belong to a finite affine closure. Its linear evaluations must separate the predictive quotient under consideration. The closure is a hypothesis, not a consequence of finite physical dimension alone. Appendix~\ref{app:quotient} gives a measurable construction and minimality statement, including the distinction between the actual quotient and the compact convex base used to certify positive mixtures.

For the audit, let $Z\in\R^k$ be bounded. There are linear maps $T_\theta:E\to\R^k$ and $c_\theta\in E^*$ such that, at predictive state $x\in D$,
\begin{equation}\label{eq:audit}
 \E_\theta[Z\mid x]=T_\theta x,\qquad
 \E_\theta[\norm Z^2\mid x]=c_\theta(x).
\end{equation}
The decoder lies in a common compact convex set $H\subset\R^k$ containing all possible conditional means. The loss is $\norm{Z-r}^2$, uniformly bounded by a declared constant $L$. All models have this same report, action, audit and loss interface. The unnormalized terminal cost is the continuous function
\begin{equation}\label{eq:terminal-cost}
 G_\theta(d,r)=c_\theta(d)-2\ip{r}{T_\theta d}+e(d)\norm r^2,
 \qquad d\in C,\quad r\in H.
\end{equation}
For a known single model, minimizing in $r$ gives the continuous perspective
\begin{equation}\label{eq:perspective}
 \overline g(d)=c(d)-\frac{\norm{Td}^2}{e(d)},\quad e(d)>0,
 \qquad \overline g(0)=0.
\end{equation}
The continuity at zero follows from compactness of $D$. In a model family the minimization in $r$ cannot be performed separately for each model.

\subsection{What the machine may retain}
A homogeneous autonomous observer has a nonempty live set $S$ and a nonempty terminal set $J$, disjoint, with
\begin{equation}\label{eq:state-budget}
 |S|+|J|\le M.
\end{equation}
Its initial live label is $s_0$. A live label $s$ selects an action according to $p_{sa}$. Having read $y$, the observer routes to $i\in S\sqcup J$ with Borel probabilities
\begin{equation}\label{eq:routing}
 \gamma_i(s,a,y)\ge0,\qquad \sum_{i\in S\sqcup J}\gamma_i(s,a,y)=1.
\end{equation}
An arrival in $J$ halts and supplies decoder $r_j\in H$. The stop/continue decision and the destination are one joint kernel. All $p,\gamma,r$ are independent of $\theta$ and of the time index. Private random bits are fresh; a retained random seed would require an additional state. The report and selected action can be processed within the current transition, but no uncharged value survives it. In particular, a terminal label is not a live label together with a free external halting bit.

Let $\tau\ge1$ be the number of acquisitions before the first terminal label. A program is admissible at budget $N$ if
\begin{equation}\label{eq:mean-budget}
 \sup_{\theta\in\Theta}\E_\theta\tau\le N<\infty.
\end{equation}
This implies almost-sure termination for each model. No statement is made about programs outside this admissible class. Define
\[
 \calR_\aut^\Theta(N,M)
 =\inf_{(p,\gamma,r)\text{ satisfying }\eqref{eq:state-budget},\eqref{eq:mean-budget}}
      \sup_{\theta\in\Theta}\E_\theta\norm{Z-r_{J_\tau}}^2.
\]
For $N\ge1,M\ge2$ the class is nonempty: acquire once and halt at a fixed decoder. For $N<1$ or $M<2$ it is empty under this interface. Its infimum is then $+\infty$, a convention about admissibility, not a divergent statistical loss.

A checkpoint with $m$ terminal labels may use the entire visible history for acquisition and stopping, subject to the same mean constraint, and then encode the stopped history into one of $m$ labels. Only this final encoding is constrained by $m$. This stronger architecture has no claimed free implementation by an $m$-state autonomous observer.

\subsection{Acquired occupation and terminal ensembles}
For one fixed program let $X_t^\theta\in D$ denote the exact predictive state after $t$ reports. Define
\begin{align}
 z_s^{\theta,\mathrm{act}}
  &=\E_\theta\sum_{t=0}^{\tau-1}\one_{\{S_t=s\}}X_t^\theta,
 &d_j^{\theta,\mathrm{act}}
  &=\E_\theta[\one_{\{J_\tau=j\}}X_\tau^\theta].\label{eq:actual-occupation}
\end{align}
These are vector integrals in the positive cone. Their masses are actual visit counts and stopping probabilities:
\[
 \sum_s e(z_s^{\theta,\mathrm{act}})=\E_\theta\tau,\qquad
 \sum_j e(d_j^{\theta,\mathrm{act}})=1.
\]
The vectors in \eqref{eq:actual-occupation} are proof objects, not registers in the observer. When solving algebraic balance equations below, a candidate $z_s^\theta$ must not be called an actual occupation until its preparation-generated part has been identified.

\section{A constructive causal acquired-geometry theorem}\label{sec:central}
We first define the finite-dimensional feasible object in the theorem. Fix live and terminal sets $S,J$ and $s_0\in S$. Put
\[
 a_s^\theta=\one_{\{s=s_0\}}x_0^\theta.
\]
A candidate consists of common action probabilities $p_{sa}$, common decoders $r_j\in H$, and, for each $\theta$, vectors $z_s^\theta,d_j^\theta\in C$ with
\begin{equation}\label{eq:budget-candidate}
 z_s^\theta-a_s^\theta\in C,\qquad
 \sum_s e(z_s^\theta)\le N,\qquad
 \sum_j e(d_j^\theta)=1.
\end{equation}
For every finite $F\subset\Theta$ and every collection $u_s^\theta,v_j^\theta\in E^*$, require
\begin{align}
 &\sum_{\theta\in F}\left[
   \sum_{s\in S}\ip{u_s^\theta}{z_s^\theta-a_s^\theta}
       +\sum_{j\in J}\ip{v_j^\theta}{d_j^\theta}\right]\notag\\
 &\quad\le\sum_{s\in S,\,a\in\calA}p_{sa}
    \int \max\left\{
      \max_{i\in S}\sum_{\theta\in F}\ip{u_i^\theta}{L_a^\theta(y)z_s^\theta},
      \max_{j\in J}\sum_{\theta\in F}\ip{v_j^\theta}{L_a^\theta(y)z_s^\theta}
              \right\}\lambda_a(dy).\label{eq:joint-support}
\end{align}
The summation over models precedes the maximum. The maximum chooses one next label, including the decision to stop, and cannot be chosen separately for different models. Define
\begin{equation}\label{eq:Phi}
 \Phi^\Theta(N,M)=
 \min\ \sup_{\theta\in\Theta}\sum_{j\in J}G_\theta(d_j^\theta,r_j),
\end{equation}
where the minimum is over the finitely many cardinalities $|S|,|J|\ge1$, $|S|+|J|\le M$, and over candidates satisfying \eqref{eq:budget-candidate}--\eqref{eq:joint-support}. The minimum and its operational meaning are part of the following result.

For the geometric part, let $K\subset\R^k$ be compact and $\nu$ a probability measure with support $K$. Write
\begin{equation}\label{eq:profile}
 \calQ_m^\nu(t)=\min_{c_1,\ldots,c_m\in\conv K}
 \ \min_{\substack{0\le h\le t\ (\nu\text{-a.e.})\\\int h\,d\nu=1}}
       \int_K\min_{1\le j\le m}\norm{x-c_j}^2h(x)\nu(dx),
 \qquad t\ge1.
\end{equation}
Both minima are attained, as proved below. The factor $t$ is a constraint on acquired probability mass, not a radius or a dimension.

\begin{definition}[Report-recoverable task with comparable acquisition]\label{def:recoverable}
A task is report-recoverable if, whenever action $a$ produces report $y$, its current conditional audit mean is a specified Borel function $x=f_a(y)\in K$, independent of the incoming history and of $\theta$, and its conditional audit variance is a common constant $B$. Suppose $\nu$ is the actual one-call law under a specified admissible reference preparation and action. There are constants $0<a_*\le1\le b_*<\infty$ such that, for every model and every positive-probability visible history, every legal next action has task-response law at most $b_*\nu$. One fixed available action $a_0$, common to all models, has conditional task-response law at least $a_*\nu$ at every such history. These are measure inequalities, not pointwise assertions on null histories.
\end{definition}
The lower inequality is needed only for the implementing action. A raw-kernel bound on a closed invariant base is one way, used in the realizations, to verify the required history quantifiers. The actual terminal law remains policy-dependent.

\begin{theorem}[Constructive budgeted causal acquired-geometry transfer]\label{thm:main}
The following statements hold.

\smallskip\noindent\textup{(i) General implementation and robust attainment.}
Under \eqref{eq:raw-instrument}--\eqref{eq:terminal-cost}, for every real $N\ge1$ and integer $M\ge2$,
\begin{equation}\label{eq:value-equality}
 \calR_\aut^\Theta(N,M)=\Phi^\Theta(N,M).
\end{equation}
A common Borel program attains this value. No uniform transience, positive stopping floor, mixing assumption, or finite bound on $\tau$ is imposed. In addition,
\begin{equation}\label{eq:finite-models}
 \calR_\aut^\Theta(N,M)
 =\sup_{\substack{F\subset\Theta\\0<|F|<\infty}}
         \calR_\aut^F(N,M).
\end{equation}
Thus every risk level strictly below the full-family optimum is already impossible on some finite model subfamily, with the same $N,M$ and one common program.

\smallskip\noindent\textup{(ii) Actual-law geometric bounds.}
For a report-recoverable task as in Definition~\ref{def:recoverable}, all fully adaptive checkpoint policies and all admissible autonomous programs satisfy
\begin{align}
 \calQ_m^\nu(b_*N)&\le\calR_\cp^\Theta(N,m)-B,
 \label{eq:profile-check-lower}\\
 \calQ_{M-1}^\nu(b_*N)&\le\calR_\aut^\Theta(N,M)-B.
 \label{eq:profile-aut-lower}
\end{align}
If $a_*N\ge1$, one common autonomous program gives
\begin{equation}\label{eq:profile-aut-upper}
 \calR_\aut^\Theta(N,M)-B
 \le\frac{b_*}{a_*}\calQ_{M-1}^\nu(a_*N).
\end{equation}
It has exactly one live state and at most $M-1$ terminal states, and its stopping time has a conditional geometric tail with mean at most $N$. The analogous checkpoint upper uses $m$ terminal states in place of $M-1$.

If every conditional response law equals $\nu$, independently of all previous reports and all actions, then for all $N\ge1$,
\begin{equation}\label{eq:exact-profile}
 \calR_\cp^\Theta(N,m)-B=\calQ_m^\nu(N),\qquad
 \calR_\aut^\Theta(N,M)-B=\calQ_{M-1}^\nu(N).
\end{equation}
The realizing rule in this case has $\E\tau=N$, including noninteger $N$.

\smallskip\noindent\textup{(iii) Matching resource law on acquired regular supports.}
Suppose in addition that, for some $s>0$, $r_0>0$ and $0<c_0\le C_0<\infty$,
\begin{equation}\label{eq:regularity}
 c_0r^s\le\nu(B(x,r))\le C_0r^s
 \quad(x\in K,\ 0<r\le r_0).
\end{equation}
Then, simultaneously for all $N\ge1$, $m\ge1$ and $M\ge2$,
\begin{align}
 \calR_\cp^\Theta(N,m)-B&\asymp (Nm)^{-2/s},\label{eq:regular-cp}\\
 \calR_\aut^\Theta(N,M)-B&\asymp \bigl(N(M-1)\bigr)^{-2/s}.\label{eq:regular-aut}
\end{align}
The constants depend only on $a_*,b_*,c_0,C_0,r_0,s$ and the diameter of $K$, not on $N,m,M$, the number of models, or the chosen adaptive policy.

\smallskip\noindent\textup{(iv) Finite construction at exactly the declared budgets.}
Under the effective raw-instrument certificate of Definition~\ref{def:effective}, for every rational $N\ge1$, integer $M\ge2$, and rational $\varepsilon>0$, a finite calculation returns a common finitely described observer and a certified risk interval of width at most $\varepsilon$. Its risk lies in the interval above $\Phi^\Theta(N,M)$, while the interval's lower endpoint bounds every admissible Borel observer. The returned observer still has mean acquisition at most $N$ and total state count at most $M$. Neither report recoverability nor posterior contraction is needed for this part. Stopping regularization and occupation-weighted error estimates, rather than a free deterministic clock or a budget relaxation, supply the construction.

\smallskip\noindent\textup{(v) Sharp-order transfer at unchanged resources.}
Consider two experiment families with the same declared interfaces and uniform preparation, joint one-call and audit errors $\delta_0,\eta,\alpha$ in the base-norm sense of Section~\ref{sec:sharp-transfer}. Put $\zeta=\delta_0+N\eta$ and $c_N=1-N^{-1}$. Their optimal values satisfy
\begin{equation}\label{eq:central-normalized}
 |\Phi_1^\Theta(N,M)-\Phi_2^\Theta(N,M)|
 \le\min\left\{L,\ L\left[2\zeta+\sqrt{\zeta^2+2c_N\zeta}\right]+\alpha\right\}.
\end{equation}
For $N=1$, $L(\delta_0+\eta)+\alpha$ is sufficient. The transfer preserves the same mean budget and total state set. Its square-root order in preparation error and in $(N-1)\eta$ is necessary in the full class, already for $M=2$, as proved by globally optimal singular-acquisition witnesses. No uniform absorption of the original competitors, normalized-filter contraction, or report recoverability is needed in this part. The bound is not asserted best possible in its numerical constants or sharp on every subclass.
\end{theorem}

Parts (i)--(iii), whose complete proofs are included, are the mean-acquisition implementation and geometric mechanisms developed in~\cite{GTFMean}. Part (iv), proved in full below, is the exact-budget constructive completion of~\cite{GTFExactBudget}. Part (v) is the budget-normalized transfer proved in Section~\ref{sec:sharp-transfer}; its all-program sharpness and the stronger memory-dependent singular frontier are proved in Section~\ref{sec:singular-memory}. Part (i) is the causal implementation theorem. Part (ii) uses a stronger report interface to turn implementation into a sharp acquired-geometry statement. Definition~\ref{def:recoverable} by itself does not require a finite affine closure; when both sets of hypotheses hold, the profile bounds are bounds on the exact value \eqref{eq:Phi}. A general normalized hidden-state filter need not satisfy recoverability. Part (iii) is not an assignment of dimension to arbitrary reachable states: its measure is an explicitly obtained reference law and its lower bound covers the actual stopped law of every competing policy.

\begin{remark}[What an occupation certificate certifies]\label{rem:ghost}
A candidate $z^\theta$ in \eqref{eq:budget-candidate} can contain invariant mass. The theorem certifies its terminal ensemble $d^\theta$, with an actual occupation no larger than $z^\theta$. It does not identify every candidate with the acquired occupation. This distinction is precisely what permits compactness without uniform absorption. Equality constraints on expected acquisition, rewards for spending resources, and negative running costs cannot be substituted in this argument.
\end{remark}

\begin{remark}[Fixed-time and stopping architectures]
The $M-1$ in \eqref{eq:exact-profile} pays for one live state. It is not the $M-n$ deduction for an internally counted, deterministic deadline. The implementing program uses no clock at all. It may run for arbitrarily many acquisitions, each charged in $\E\tau$. A comparison with a pathwise cap, a latency quantile, or a prescribed audit time requires a different resource constraint.
\end{remark}

\section{Common routing and removal of invariant occupation}\label{sec:occupation-proof}
We prove Theorem~\ref{thm:main}(i) without using any realization or any metric geometry.

\begin{lemma}[A joint routing separation criterion]\label{lem:routing}
Fix common $p$ and vectors $z_s^\theta\in C$. There exist common Borel routing probabilities \eqref{eq:routing} such that, for all $\theta$,
\begin{align}
 z_i^\theta-a_i^\theta
  &=\sum_{s,a}p_{sa}\int\gamma_i(s,a,y)L_a^\theta(y)z_s^\theta\,\lambda_a(dy)
       &&(i\in S),\label{eq:balance-live}\\
 d_j^\theta
  &=\sum_{s,a}p_{sa}\int\gamma_j(s,a,y)L_a^\theta(y)z_s^\theta\,\lambda_a(dy)
       &&(j\in J)\label{eq:balance-terminal}
\end{align}
if and only if \eqref{eq:joint-support} holds for every finite $F$ and every displayed dual vector.
\end{lemma}
\begin{proof}
For each pair $(s,a)$ the simplex-valued kernels form a weak-star compact convex subset of $L^\infty(\lambda_a)^{S\sqcup J}$. Positivity and the sum-one constraint are weak-star closed: test them against nonnegative $L^1$ functions. The product over finitely many $(s,a)$ is compact and convex. Integration against $L_a^\theta(y)z_s^\theta$ is weak-star continuous by \eqref{eq:raw-instrument}. For a finite $F$, the vector of the right sides of \eqref{eq:balance-live}--\eqref{eq:balance-terminal} therefore has a compact convex image.

Its support function is the right side of \eqref{eq:joint-support}. Indeed, after summing the paired contributions over $\theta\in F$, one maximizes a linear function over a probability simplex, for each $(s,a,y)$. A maximizing label can be selected by the smallest-index rule among finitely many measurable functions. This gives a Borel maximizer and the stated integral of a maximum. Finite-dimensional separation proves that the target vector belongs to the image exactly when all support inequalities hold.

For an arbitrary $\Theta$, impose the balance equations first on each finite $F$. Their solution sets in the same weak-star compact routing space are closed and have the finite intersection property: the union of finitely many finite model sets is finite. A common routing kernel exists by compactness. On standard Borel report spaces its equivalence classes admit Borel versions. Changing representatives on a common $\lambda_a$-null set and assigning a fixed destination there makes \eqref{eq:routing} pointwise valid without changing any model law. Necessity is integration of the pointwise simplex bound.
\end{proof}

For such a kernel define positive linear maps on ensembles,
\[
 (P_\theta z)_i=\sum_{s,a}p_{sa}\int\gamma_i(s,a,y)L_a^\theta(y)z_s\,\lambda_a(dy),
 \quad
 (Q_\theta z)_j=\sum_{s,a}p_{sa}\int\gamma_j(s,a,y)L_a^\theta(y)z_s\,\lambda_a(dy).
\]
Write $\boldsymbol e(z)=\sum_s e(z_s)$, with the analogous convention on terminal ensembles. Normalization gives
\begin{equation}\label{eq:mass-conservation}
 \boldsymbol e(P_\theta z)+\boldsymbol e(Q_\theta z)=\boldsymbol e(z)
 \quad\text{for every positive live ensemble }z.
\end{equation}

\begin{lemma}[Budget-preserving removal]\label{lem:removal}
Let $P,Q$ be positive ensemble maps satisfying \eqref{eq:mass-conservation}. If a finite positive ensemble solves $z=a+Pz$, then
\begin{equation}\label{eq:removal}
 h=\lim_{n\to\infty}P^nz\in C^S,
 \qquad Ph=h,\qquad Qh=0,
 \qquad z^{\mathrm{act}}=\sum_{n=0}^\infty P^na=z-h.
\end{equation}
The common observer started at preparation $a$ has actual continuing ensemble $z^{\mathrm{act}}$, terminal ensemble $Qz$, and mean acquisition $\boldsymbol e(z^{\mathrm{act}})\le\boldsymbol e(z)$. Every positive linear running-cost functional is no larger on the actual ensemble.
\end{lemma}
\begin{proof}
Since $Pz=z-a\le_C z$, the sequence $P^nz$ decreases in the cone order. Its scalar mass decreases to a finite limit. Compactness of the base implies an estimate $\norm w\le C_D\boldsymbol e(w)$ for positive ensembles. Thus, for $n\ge m$,
\[
 \norm{P^mz-P^nz}\le C_D\bigl(\boldsymbol e(P^mz)-\boldsymbol e(P^nz)\bigr),
\]
so the sequence converges in the finite-dimensional space. Its limit $h$ is positive, and continuity of $P$ gives $Ph=h$. Applying \eqref{eq:mass-conservation} to $h$ gives $\boldsymbol e(Qh)=0$. Strict positivity of the mass functional on the cone implies $Qh=0$, component by component.

Iteration of $z=a+Pz$ yields
\[
 z=\sum_{n=0}^{t-1}P^na+P^tz.
\]
Taking the limit proves \eqref{eq:removal}. Starting the physical experiment at $a$, induction with the raw instruments shows that $P^na$ is exactly the unnormalized predictive ensemble of surviving live labels after $n$ acquisitions. This statement uses only linearity before normalization and the fact that the same report-routing kernel is executed at every step. Tonelli's theorem, applied first to masses and then to the bounded base coordinates, gives the actual occupation and $\E\tau=\boldsymbol e(z-h)<\infty$. Hence $\tau<\infty$ almost surely. Its terminal ensemble is $Q\sum_nP^na=Q(z-h)=Qz$. A positive linear cost decreases when the positive component $h$ is removed.
\end{proof}

\begin{example}[Why the candidate is not the acquired law]\label{ex:phantom}
Take the scalar cone $\R_+$, live states $0,1$, terminal state $2$, preparation $a=(1,0)$. State $0$ always halts after one call; state $1$ always returns to itself. Then
\[
 P=\begin{pmatrix}0&0\\0&1\end{pmatrix},\qquad Q=(1,0),\qquad z=(1,h),\quad d=1
\]
solves the balance equations for every $h\ge0$. The prepared run never visits state $1$. Its mean acquisition is $1$, not $1+h$. The extra component changes neither the terminal law nor an upper budget constraint. This is the elementary phantom-measure phenomenon in occupation theory~\cite{DPR}; its explicit role here is to delimit the interpretation of a predictive certificate.
\end{example}

\begin{proof}[Proof of Theorem~\ref{thm:main}(i)]
For an admissible observer, linear conditional integration and Tonelli give \eqref{eq:balance-live}--\eqref{eq:balance-terminal} with the actual ensembles \eqref{eq:actual-occupation}. The mean budget gives \eqref{eq:budget-candidate}; Lemma~\ref{lem:routing} gives \eqref{eq:joint-support}. Conditional expectation at the protected audit gives risk $\sum_jG_\theta(d_j^\theta,r_j)$. Therefore $\Phi^\Theta\le\calR_\aut^\Theta$.

Conversely, a feasible candidate has one common routing witness by Lemma~\ref{lem:routing}. For each $\theta$, Lemma~\ref{lem:removal} identifies the actual ensemble generated by its preparation. The same kernel, not a model-specific choice of kernel, is used for every $\theta$. Its actual expected acquisition is at most $N$, and its terminal ensemble is exactly $d^\theta$. Its loss is therefore the candidate loss. This proves the reverse inequality, including at the level of infima.

It remains to show attainment. For fixed cardinalities, action rows lie in compact simplices and decoders in $H^J$. For each model, the positive ensemble sets in \eqref{eq:budget-candidate} are compact, because their masses are bounded and $D$ is compact. Take the product of these compact sets over $\Theta$. Each inequality \eqref{eq:joint-support} depends on finitely many models and is continuous in $p,z,d$. To verify continuity of its integral, bound the difference of maxima by the largest difference of the finitely many linear terms, and use the integrability in \eqref{eq:raw-instrument}. The feasible set is therefore compact. It is nonempty by the one-call stopping program. The objective, a supremum of continuous model losses, is lower semicontinuous; it attains its minimum. There are only finitely many cardinality choices. The implementing kernel supplied above attains the resulting global minimum.

For \eqref{eq:finite-models}, the inequality from right to left follows by restriction. Suppose that $r$ is strictly smaller than the full-family optimum and no finite subfamily has optimum greater than $r$. Work in the finite disjoint union of the compact candidate spaces, adding the closed constraints $\sum_jG_\theta(d_j^\theta,r_j)\le r$. Any finite collection of support inequalities and risk constraints refers to finitely many models. A minimizing candidate for their union satisfies the collection. Coordinates for other models may be filled arbitrarily in their nonempty budget sets. Thus these closed sets have the finite intersection property and yield a full-family feasible candidate of risk at most $r$, a contradiction. This also proves the strict finite-model obstruction. There is no interchange of $\inf$ and $\sup$, and no bound on the size of the obstructing finite set is asserted.
\end{proof}

\begin{corollary}[Additional nonnegative costs]\label{cor:costs}
In Theorem~\ref{thm:main}(i), finitely many further constraints of the form
\[
 \sum_{s,a}p_{sa}\,\kappa_{\ell,sa}^\theta(z_s^\theta)\le B_\ell,
\]
where each $\kappa_{\ell,sa}^\theta$ is a fixed nonnegative linear functional representing an actual conditional expected primitive cost, can be added. The unit acquisition bound is retained. The value equality and finite-model statement remain valid, and the value is attained whenever the constrained admissible class is nonempty.
\end{corollary}
\begin{proof}
The additional inequalities are closed. After extracting a common routing witness, removal of its invariant occupation can only decrease each displayed cost. The compactness and finite-intersection arguments are unchanged.
\end{proof}

\begin{remark}[Necessity versus a finite computation]
The support inequalities are an exact realizability criterion, not a finite numerical test in general. Infeasibility for a fixed candidate and finite model family has a strict finite-dimensional separating vector. The global finite-model conclusion is stronger but still non-effective: it supplies neither a uniformly bounded witness size nor an algorithm for locating it. These distinctions prevent a compactness proof from being reported as a finite-bit synthesis theorem.
\end{remark}

\section{Exact-budget construction and occupation-weighted stability}\label{sec:effective}
The realization criterion does not by itself make its Borel witness computable. This section closes that gap under effective hypotheses on the unnormalized experiment, not on an optimal observer or a normalized filter. Three different operations must be distinguished: stopping regularization creates budget slack, approximation perturbs an executable kernel, and finite evaluation certifies the perturbed program. None changes the persistent-state set.

\subsection{Stopping regularization without an extra label}
For a fixed program write $P$ for its live-ensemble map and $Q$ for its terminal map, as in Section~\ref{sec:occupation-proof}. Choose an already available terminal label $j_0$. For $0<\rho\le1$, first sample the original routing decision. If it is live, replace it by $j_0$ with independent probability $\rho$; otherwise leave it unchanged. The physical state after the call is not reset. In ensemble notation,
\begin{equation}\label{eq:killed-routing}
 P_\rho=(1-\rho)P,\qquad Q_\rho=Q+\rho RP,
\end{equation}
where $R$ sends every live component, with its same physical vector, to terminal component $j_0$. This is a homogeneous report-only program. Its conditional stopping probability at every possible live configuration is at least $\rho$.

\begin{lemma}[A paid mean margin]\label{lem:margin}
For every admissible program, every model, and $0<\rho\le1$, the regularized program uses the same labels and satisfies
\begin{align}
 \E_\theta\tau_\rho
 &\le1+(1-\rho)(\E_\theta\tau-1),\label{eq:mean-margin}\\
 |R_\theta(\gamma_\rho)-R_\theta(\gamma)|
 &\le L\rho(\E_\theta\tau-1).\label{eq:killing-risk}
\end{align}
Consequently, for $1\le N_1\le N_2$ and fixed $M\ge2$,
\begin{equation}\label{eq:budget-lipschitz}
 0\le\calR_\aut^\Theta(N_1,M)-\calR_\aut^\Theta(N_2,M)
       \le L(N_2-N_1).
\end{equation}
This statement requires only a common bounded task and the availability of the stated terminal fallback; no finite predictive closure is needed.
\end{lemma}
\begin{proof}
Couple the original experiment and the modified one until the latter stops. For the unmodified stopping time $\tau$ and independent geometric $G$ with success probability $\rho$, the modified stopping time is $\min(\tau,G)$. For $t\ge1$,
\[
 \Pp(\tau_\rho>t)=(1-\rho)^t\Pp(\tau>t).
\]
Summing the tails and bounding $(1-\rho)^t\le1-\rho$ proves \eqref{eq:mean-margin}. The programs and their terminal audits agree unless a live routing decision is replaced. The union bound gives
\[
 \Pp(G<\tau)\le\rho\sum_{t\ge1}\Pp(\tau>t)
                  =\rho(\E\tau-1).
\]
Both losses belong to $[0,L]$, so this proves \eqref{eq:killing-risk}, including a possible change in the physical state at the earlier audit. For $N_2>1$, set $\rho=(N_2-N_1)/(N_2-1)$, using the unchanged program when $\rho=0$. Apply the two inequalities uniformly over models and then take infima. The case $N_2=1$ is immediate. No mixture seed needs to persist: the replacement coin is fresh on every live transition.
\end{proof}

\subsection{Two norms for two different obligations}
On $E$ use the base norm
\[
 \bnorm v=\inf\{e(v_+)+e(v_-):v=v_+-v_-,\ v_\pm\in C\},
\]
and on an ensemble its sum over labels. A positive mass-preserving map is a contraction in these norms; this follows by applying it to each positive part and taking the infimum. In particular, the map $(P,Q)$ from a live ensemble to its live and terminal successors is contractive. For a regularized program, $\bnorm{P_\rho}\le1-\rho$ and $\bnorm{(I-P_\rho)^{-1}}\le\rho^{-1}$.

Consider two implementations on the same live and terminal labels and the same physical ordered space. Denote their maps and preparations by $(P,Q,a)$ and $(\widetilde P,\widetilde Q,\widetilde a)$. The first is admissible and has actual occupation $z=\sum_{t\ge0}P^ta$ and mean $T=\boldsymbol e(z)$. The second has a stopping floor $\rho$. Define its occupation-weighted defect against the first by
\begin{equation}\label{eq:weighted-defect}
 D_z=\bnorm{\widetilde a-a}
       +\bnorm{(\widetilde P-P)z}
       +\bnorm{(\widetilde Q-Q)z}.
\end{equation}
This is measured on the preparation-generated occupation, not an arbitrary solution of a balance equation.

\begin{lemma}[Terminal risk versus expected acquisition]\label{lem:resolvent-error}
Under these hypotheses both terminal laws exist, and
\begin{align}
 \bnorm{\widetilde d-d}&\le D_z,\label{eq:terminal-defect}\\
 |\E\widetilde\tau-\E\tau|&\le D_z/\rho.\label{eq:mean-defect}
\end{align}
If the same decoder is used and the terminal cost functionals differ by at most $\alpha$ per unit positive mass, then
\begin{equation}\label{eq:risk-defect}
 |\widetilde R-R|\le LD_z+\alpha.
\end{equation}
In particular, if $\bnorm{\widetilde a-a}\le\delta_0$ and the joint one-call maps differ in induced base norm by at most $\eta$, then $D_z\le\delta_0+\eta T$. No stability of the normalized posterior map is assumed.
\end{lemma}
\begin{proof}
Write $\widetilde Z=(I-\widetilde P)^{-1}$ and $\widetilde K=\widetilde Q\widetilde Z$. Uniform killing implies convergence of the series defining $\widetilde Z$. The positive map $\widetilde K$ preserves live mass: summing conservation through $t$ transitions leaves a surviving mass at most $(1-\rho)^t$ times the initial mass. Hence $\bnorm{\widetilde K}\le1$. Since $z=a+Pz$ and $\widetilde z=\widetilde Z\widetilde a$,
\[
 \widetilde z-z=\widetilde Z\bigl[(\widetilde a-a)+(\widetilde P-P)z\bigr].
\]
The norm bound on $\widetilde Z$ and the norm-one mass functional prove \eqref{eq:mean-defect}. Moreover,
\[
 \widetilde d-d=\widetilde K\bigl[(\widetilde a-a)+(\widetilde P-P)z\bigr]
                         +(\widetilde Q-Q)z.
\]
The map $(u,v)\mapsto\widetilde Ku+v$ is contractive for the sum base norm. This proves \eqref{eq:terminal-defect}. A terminal cost between $0$ and $L$ on every normalized positive component has dual base norm at most $L$. Applying it, and then accounting for its change on the unit terminal mass, proves \eqref{eq:risk-defect}. Finally $z$ is positive and $\bnorm z=T$, giving the last assertion. The proof uses unnormalized finite-dimensional positive maps, not a joint classical hidden-state interpretation of a quantum experiment.
\end{proof}

The factors in \eqref{eq:terminal-defect} and \eqref{eq:mean-defect} have different meanings. A rare routing discrepancy has small probability but may be followed by a long surviving excursion. For a one-live-state program whose stopping probability changes from $\rho$ to $\rho+u$, the mean changes by $u/[\rho(\rho+u)]$. Thus a mean-cost estimate cannot simply be replaced by the bounded-risk estimate.

\begin{proposition}[Same-budget stability of the optimal risk]\label{prop:same-budget}
Suppose two model families have the same labels, actions, report interface, ordered physical space, and bounded terminal task. Uniformly over models and all common routing programs, assume preparation error at most $\delta_0$, joint one-call error at most $\eta$, and terminal cost error at most $\alpha$. Put $\zeta=\delta_0+N\eta$. For $N\ge1$ and $M\ge2$,
\begin{equation}\label{eq:optimal-modulus}
 \left|\calR_{\aut,1}^\Theta(N,M)-\calR_{\aut,2}^\Theta(N,M)\right|
 \le \min\left\{L,
 L\left[\sqrt{2(N-1)\zeta}+\zeta\right]+\alpha\right\}.
\end{equation}
The program transferred in either direction respects the same $N,M$. The exponent in this sufficient modulus is not asserted optimal.
\end{proposition}
\begin{proof}
First let $N>1$ and $0<\zeta<(N-1)/2$. Apply Lemma~\ref{lem:margin} to an arbitrary admissible program with
$\rho=\sqrt{2\zeta/(N-1)}$. Its mean in the first experiment is at most $N-\rho(N-1)$. In the second, the same regularized routing still has stopping floor $\rho$. Lemma~\ref{lem:resolvent-error} bounds its mean by
\[
 N-\rho(N-1)+\zeta/\rho
       =N-\tfrac12\rho(N-1)\le N.
\]
Its risk has increased by at most $L\rho(N-1)+L\zeta+\alpha$. Take infima, and exchange the two experiments for the reverse inequality. If $\zeta\ge(N-1)/2$, set $\rho=1$. Both programs then halt after one call, their means are exactly one, and the same argument gives $L(N-1)+L(\delta_0+\eta)+\alpha$, no greater than the second quantity in \eqref{eq:optimal-modulus}. If $N=1$, every admissible run stops on the first call; send all originally live routes to $j_0$, which changes no actual law, and use the one-call estimate. If $\zeta=0$, the actual maps and costs agree to their stated tolerance, and only $\alpha$ remains. Since both values lie in $[0,L]$, the truncation by $L$ is valid.
\end{proof}

\subsection{Effective primitive hypotheses}
The next definition is additional input for finite synthesis. It is not part of the geometric or occupation identities.

\begin{definition}[Effective raw-instrument certificate]\label{def:effective}
In addition to Section~\ref{sec:experiments}, suppose the following data are given.
\begin{enumerate}[label=\textup{(E\arabic*)}]
\item The model set is an effectively compact metric space: finite nets with a certified covering radius can be constructed. The preparation, the action instruments in the integral induced base norm, and the terminal cost functionals have computable uniform moduli in the model parameter. Their values on the finite nets are computable. The physical dimension and cone norm are fixed and computably represented.
\item For every $h>0$ a finite common report partition $\mathcal B_{a,h}$ is supplied. For positive-reference-mass cells $B$, put
\[
 \overline L_{a,h}^\theta(y)=\lambda_a(B)^{-1}\int_B L_a^\theta(v)\lambda_a(dv)
 \quad(y\in B).
\]
There is a computable bound $\eta_h\downarrow0$ with
\begin{equation}\label{eq:report-approximation}
 \sup_{\theta,a}\int\bnorm{L_a^\theta(y)-\overline L_{a,h}^\theta(y)}\lambda_a(dy)
       \le\eta_h.
\end{equation}
The cell instrument integrals can be computed with certified error. Reference-null cells may be assigned any fixed route. The executed partition selector is either exact under the declared digitization primitive or has a uniform, certified induced-instrument defect included in $\eta_h$.
\item The common decoder body has finite computable legal nets with uniformly controlled loss error. In particular, at a decoder mesh $s$ there is a finite list of finitely described legal outputs whose replacement changes every terminal cost by at most $\omega_H(s)\downarrow0$. The lists, not a model-dependent decoder, are part of the program.
\item Computable positive-map and audit integrals, together with the preceding norm bounds, permit interval evaluation of a finite product to arbitrary absolute error. The running acquisition cost is exactly one per call. $N$ is rational when a finite algorithm is requested.
\end{enumerate}
\end{definition}
These assumptions include the cost of obtaining the numerical report partition. They do not turn an arbitrary Borel membership problem into a finite operation. They are verified below from classical and quantum raw instruments, including examples with arbitrarily slow or absent physical mixing.

\begin{lemma}[Uniform reduction of every Borel program]\label{lem:finite-routing}
Fix a partition as in \eqref{eq:report-approximation}, a decoder net, an integer $q\ge1$, and a dyadic $0<\rho\le1$. Every regularized Borel program on $S,J$ has a program on the same labels, using only the current report cell and probabilities with dyadic denominator $Q=2^q$ before regularization. Both have stopping floor $\rho$. Their joint one-call maps differ by at most
\begin{equation}\label{eq:eta-compile}
 \eta=\eta_h+\frac{2(|\calA|-1)}{Q}
                      +\frac{2(|S|+|J|-1)}{Q}.
\end{equation}
The decoder replacement costs at most $\omega_H(s)$. The finite grid used here contains approximations to all legal Borel programs, not only smooth or previously constructed observers.
\end{lemma}
\begin{proof}
For the unregularized routing average $\gamma_i(s,a,y)$ over each cell with respect to the common reference measure. This produces common simplex entries $\beta_i(s,a,B)$, independent of the true model. The finite-cell program executes the actual cell instrument, not a conditional experiment. Its $i$th output is the integral of $\gamma_i\overline L$; hence the difference from the original is the integral of $\gamma_i(L-\overline L)$. Summing norms over destinations, using nonnegativity and $\sum_i\gamma_i=1$, gives at most $\eta_h$ for each input physical unit state and live label. The same bound for signed inputs follows from the base-norm decomposition.

A $d$-simplex has a $Q$-grid point within $2(d-1)/Q$ in $\ell^1$: round its first $d-1$ coordinates down and give the remainder to the last. Apply this first to action rows and then to each report-routing row. Normalization and positivity of the physical instrument make their errors at most the corresponding $\ell^1$ errors in the joint ensemble norm. Triangle inequality gives \eqref{eq:eta-compile}. Averaging and rounding the original rows, then applying the same regularization, is legitimate because \eqref{eq:killed-routing} is a positive mass-preserving postprocessing of the joint next-label output. It cannot enlarge the error. It also leaves a stopping floor $\rho$ regardless of rounding. Finally use the supplied legal decoder net.
\end{proof}

\begin{theorem}[Finite certificates at the same mean and state budgets]\label{thm:effective-budget}
Under Definition~\ref{def:effective}, for every rational $N\ge1$, integer $M\ge2$, and rational accuracy $\varepsilon>0$, a finite algorithm returns a common finitely described observer $\widehat\gamma$ and rational numbers $\ell,u$ such that
\begin{equation}\label{eq:exact-budget-certificate}
 \sup_\theta\E_\theta\widehat\tau\le N,\qquad
 |\widehat S|+|\widehat J|\le M,\qquad
 \ell\le\Phi^\Theta(N,M)=\calR_\aut^\Theta(N,M)
       \le\sup_\theta R_\theta(\widehat\gamma)\le u,
 \qquad u-\ell\le\varepsilon.
\end{equation}
No stopping floor is imposed on competing observers, and no uniform absorption assumption is imposed on the original experiment. The construction is finite, not claimed computationally efficient.
\end{theorem}
\begin{proof}
Assume first $N>1$. Choose a positive dyadic $\rho<1$ sufficiently small, and put $\sigma=\rho(N-1)$. Choose the common report partition and probability grid so fine that
\begin{equation}\label{eq:reserve-choice}
 \eta N/\rho\le\sigma/8.
\end{equation}
Choose the decoder net sufficiently fine. Enumerate every cardinality pair $|S|,|J|\ge1$, $|S|+|J|\le M$, every grid action and routing table, and every tuple of net decoders; use a fixed first live label initially. Apply the same regularization to every table. This is a finite class independent of the model.

We describe a certified test, without assuming that feasibility at a boundary is decidable. Every candidate has stopping floor $\rho$, so its mean is at most $1/\rho$ from any normalized live ensemble. For two nearby models with preparation difference $d_0$, joint instrument difference $d_1$, and terminal cost difference $d_f$, Lemma~\ref{lem:resolvent-error} gives uniform moduli
\begin{equation}\label{eq:evaluation-moduli}
 \Delta_T\le d_0/\rho+d_1/\rho^2,
 \qquad \Delta_R\le L(d_0+d_1/\rho)+d_f.
\end{equation}
For an inexact executed selector with induced-instrument defect $\eta_{\mathrm{sel}}$, include the additional mean bound $\eta_{\mathrm{sel}}/\rho^2$ and risk bound $L\eta_{\mathrm{sel}}/\rho$ in the evaluation upper intervals. Choose its accuracy along with the partition so these bounds fit the tolerances below; the physical executed table, rather than only its ideal cell surrogate, is being certified. Use a model net and certified finite-product evaluation to enclose the suprema of the candidate's mean and risk. Truncating the live series after $k$ terms leaves mean at most $(1-\rho)^k/\rho$, and truncating terminal accumulation leaves risk at most $L(1-\rho)^k$. Thus finite $k$, finite model nets, and sufficiently small arithmetic intervals yield upper bounds $T^+,R^+$ with
\begin{equation}\label{eq:eval-errors}
 0\le T^+-\sup_\theta\E_\theta\tau\le\sigma/4,
 \qquad 0\le R^+-\sup_\theta R_\theta\le\epsilon_e.
\end{equation}
Accept only tables for which the computed rational $T^+\le N$. Unknown or boundary cases may be discarded; no false feasibility is accepted.

For any admissible original program, Lemma~\ref{lem:margin} first gives a program with mean at most $N-\sigma$ and risk increase at most $L\sigma$. Lemma~\ref{lem:finite-routing} and the occupation-sensitive version of Lemma~\ref{lem:resolvent-error}, with the original mean at most $N$, yield a grid program satisfying
\[
 \sup_\theta\E_\theta\tau_{\mathrm{grid}}\le N-7\sigma/8,
 \qquad \sup_\theta R_\theta(\gamma_{\mathrm{grid}})
       \le\sup_\theta R_\theta(\gamma)+L\sigma+L\eta N+\omega_H(s).
\]
By \eqref{eq:eval-errors} this grid program is accepted. In particular the finite accepted set is nonempty, and this covering statement holds for every competing Borel program, including a minimizing program supplied by Theorem~\ref{thm:main}(i).

Let $u$ be the smallest accepted $R^+$ and choose its program. Since every accepted program is truly feasible, $\calR_\aut^\Theta\le u$. Conversely the preceding approximation and \eqref{eq:eval-errors} give
\[
 u\le\calR_\aut^\Theta(N,M)+E,
 \qquad E=L\sigma+L\eta N+\omega_H(s)+\epsilon_e.
\]
The parameters can be chosen effectively so that a rational upper bound $E^+\le\varepsilon$ holds. Set $\ell=\max\{0,u-E^+\}$, increasing $u$ only by an already included rational rounding allowance. Then \eqref{eq:exact-budget-certificate} follows. The lower bound covers all Borel programs because every such feasible program has an accepted approximation; it is not a lower bound inferred from an unsuccessful local search.

For $N=1$, every admissible run stops after its first report under every model. Replacing all live destinations of its first action row by $j_0$ therefore changes no actual law. It has a one-live-label implementation with at most the same number of terminal labels. Enumerate only one-call finite tables with no live routing. Their mean is identically one, so no budget comparison or margin is needed. The report, probability, model and decoder approximations give the same two-sided certificate by a single integral and finite enumeration. This treats the endpoint without taking a nonuniform limit in $\sigma$.
\end{proof}

Theorem~\ref{thm:effective-budget} proves part (iv) of Theorem~\ref{thm:main}. Its finite enumeration controls the global risk value; it does not certify a general closed-form dimension law or a polynomial search complexity. The resource--resolution rates in the next section require actual acquisition geometry in addition to the effective raw data.

\section{Budget-normalized causal risk transfer}\label{sec:sharp-transfer}
A stopped task has two different sensitivities. Terminal risk is measured on a probability law, whereas feasibility is measured on the entire live occupation. The preceding resolvent argument respects this distinction. We now use the full mean constraint, rather than its first linear consequence, to obtain a bound with the correct dependence on the acquisition budget. The construction still changes only the routing to an existing terminal label.

\subsection{The exact survival envelope}
For $N\ge1$ put $n=\lfloor N\rfloor$ and $b=N-n$. For $0<\rho\le1$ and $q=1-\rho$, define
\begin{align}
 H_N(\rho)&=\sum_{t=0}^{n-1}q^t+bq^n,
 &\Delta_N(\rho)&=N-H_N(\rho),\label{eq:survival-envelope}\\
 J_N(\rho)&=1-q^{n-1}(1-b\rho).
 \label{eq:replacement-envelope}
\end{align}
The convention $0^0=1$ is used. The temporary symbols $n,b,q$ in this subsection do not denote a deterministic acquisition horizon, a stratum mass, or a likelihood ratio.

\begin{lemma}[Sharp envelopes under a mean constraint]\label{lem:sharp-survival}
Let $\tau$ take values in $\{1,2,\ldots\}$ with $\E\tau\le N$, and let $G$ be independent geometric with success probability $\rho$ on the same set. Then
\begin{equation}\label{eq:sharp-survival-bounds}
 \E\min(\tau,G)\le H_N(\rho),\qquad
 \Pp(G<\tau)\le J_N(\rho)\le\min\{1,\rho(N-1)\}.
\end{equation}
Both first bounds are attained, among all such laws of $\tau$, by putting probabilities $1-b,b$ on $n,n+1$. In addition,
\begin{equation}\label{eq:quadratic-margin}
 \Delta_N(\rho)\ge
 \frac{\rho N(N-1)}{2(1+\rho N)}.
\end{equation}
This last lower bound is not asserted exact.
\end{lemma}
\begin{proof}
For an integer $k\ge1$ the conditional mean of $\min(k,G)$ is
$f(k)=\sum_{t=0}^{k-1}q^t$, and the replacement probability is
$g(k)=1-q^{k-1}$. The successive increments of both functions are nonnegative and nonincreasing. Their piecewise linear interpolants on $[1,\infty)$ are therefore increasing and concave. Jensen's inequality, followed by $\E\tau\le N$, proves the first two assertions, including their equality case. The union bound $1-q^{k-1}\le\rho(k-1)$ proves the remaining bound in \eqref{eq:sharp-survival-bounds}.

For every integer $t\ge0$,
$q^t\le(1+t\rho)^{-1}$: for $0<\rho<1$ this follows from
$(1-\rho)^{-t}\ge1+t\rho/(1-\rho)\ge1+t\rho$, and the endpoints follow directly. Consequently
\[
 \Delta_N(\rho)=\sum_{t=0}^{n-1}(1-q^t)+b(1-q^n)
 \ge \frac{\rho}{1+\rho N}
       \left[\frac{n(n-1)}2+bn\right].
\]
The bracket equals $N(N-1)/2+b(1-b)/2$, proving
\eqref{eq:quadratic-margin}. The equality law in the first assertion is a probability witness, not an uncharged clock in the observer.
\end{proof}

\subsection{A transfer at unchanged resources}
Write $\calR_i^\Theta(N,M)$ for the autonomous value in experiment family $i=1,2$. The two families have the same finite ordered physical space, actions, report reference spaces, legal decoder body and bounded terminal task interface. They may have different preparations, instruments and terminal cost functionals. For each common routing program suppose, uniformly in $\theta$, that the preparation error in base norm is at most $\delta_0$, the induced base-norm error of the joint one-call live/terminal map is at most $\eta$, and the terminal cost functional error per unit positive mass is at most $\alpha$. Integral raw-instrument bounds imply the one-call condition by positivity; no normalized posterior distance is assumed. Every loss belongs to $[0,L]$ in both families.

\begin{theorem}[Sharp-order same-budget transfer]\label{thm:normalized-transfer}
For every real $N\ge1$ and integer $M\ge2$, put
$\zeta=\delta_0+N\eta$ and $c_N=1-N^{-1}$. Then
\begin{equation}\label{eq:normalized-transfer}
 \left|\calR_1^\Theta(N,M)-\calR_2^\Theta(N,M)\right|
 \le \min\left\{L,\ L\left[2\zeta+
          \sqrt{\zeta^2+2c_N\zeta}\right]+\alpha\right\}.
\end{equation}
At $N=1$ the stronger bound $\min\{L,L(\delta_0+\eta)+\alpha\}$ holds. The transfer is effected in each direction by a parameter-independent program with the same state set and the same mean budget. No uniform absorption of the original competitors is required.

The square-root order in $\delta_0$, and in $(N-1)\eta$ for $N\ge2$, is attained up to numerical constants within this experiment class. The witnesses in Theorems~\ref{thm:rare-preparation} and~\ref{thm:rare-instrument} minimize over all two-state programs, not just a chosen regularization family. The constants in \eqref{eq:normalized-transfer} are not claimed best possible.
\end{theorem}
\begin{proof}
Fix any admissible source program. Apply the existing-label regularization of \eqref{eq:killed-routing}. The coupling in Lemma~\ref{lem:margin} and Lemma~\ref{lem:sharp-survival} give source mean at most $H_N(\rho)$ and risk increase at most $LJ_N(\rho)$. The identical regularized program in the other family has stopping floor $\rho$. By Lemma~\ref{lem:resolvent-error}, its mean is at most
\begin{equation}\label{eq:normalized-feasibility}
 H_N(\rho)+\frac{\delta_0+\eta H_N(\rho)}{\rho}
 \le H_N(\rho)+\frac{\zeta}{\rho},
\end{equation}
and its additional risk error is at most $L\zeta+\alpha$. Thus
$\rho\Delta_N(\rho)\ge\zeta$ is a sufficient feasibility certificate.

Assume $N>1$ and $\zeta>0$. Put
\[
 s=\frac{N\zeta}{N-1},\qquad t=s+\sqrt{s^2+2s}.
\]
If $t\le N$, choose $\rho=t/N$. Since $t^2=2s(1+t)$,
\eqref{eq:quadratic-margin} gives $\rho\Delta_N(\rho)\ge\zeta$.
Furthermore
\[
 \rho(N-1)=\zeta+\sqrt{\zeta^2+2c_N\zeta}.
\]
Equations~\eqref{eq:sharp-survival-bounds} and
\eqref{eq:normalized-feasibility} prove the stated one-sided transfer.

If $t>N$, choose $\rho=1$. Both regularized programs stop after exactly one call, so feasibility is exact, independently of any norm estimate. The risk increase is at most $L(N-1)+L(\delta_0+\eta)+\alpha$. The condition $t>N$ is equivalent to
\[
 \zeta>\frac{N(N-1)}{2(N+1)}.
\]
At equality one has $\zeta+\sqrt{\zeta^2+2c_N\zeta}=N-1$; this expression is increasing in $\zeta$. Since $\delta_0+\eta\le\zeta$, the same displayed bound follows.

For $N=1$, every admissible source run already halts after one call. Reroute all formally live outcomes to the fixed terminal fallback. This changes no source execution law, and makes the target mean exactly one; use the one-call norm bound. For $\zeta=0$, the only remaining error is the stated audit error. Take infima over source programs and then exchange the families. Both values lie in $[0,L]$, which justifies the truncation.
\end{proof}

The sharper envelope itself is also an executable certificate. For any $\rho$ satisfying
\begin{equation}\label{eq:exact-envelope-certificate}
 \rho\Delta_N(\rho)\ge\delta_0+\eta H_N(\rho),
\end{equation}
the proof gives error at most
$L[J_N(\rho)+\delta_0+\eta H_N(\rho)]+\alpha$.
A one-call fallback is available even if this sufficient inequality has no solution. In particular, the perturbation enters through the actual paid scale $\delta_0+N\eta$, not through an additional factor $N$ multiplying that scale. The result concerns reoptimized constrained values. It does not say that the unregularized occupation map is continuous, or that one unchanged boundary-optimal program remains feasible after perturbation. This distinction is essential in absorbing control~\cite{PZ,DPR}.

\begin{corollary}[A sufficient finite-precision allocation]\label{cor:normalized-precision}
Suppose $N\ge2$ and $0<a\le1/4$. For any admissible source program choose $\rho=a/N$. If an implementation on the same labels preserves this stopping floor, has preparation and one-call defects satisfying
\begin{equation}\label{eq:precision-allocation}
 \delta_0+N\eta\le a^2/64,
\end{equation}
and terminal cost error at most $\alpha$, it is feasible at the same $N,M$ and has risk at most the source risk plus
$L(a+a^2/64)+\alpha$.
In particular, one-call error of order $a^2/N$ is sufficient, rather than order $a^2/N^2$. This is a sufficient precision region, not an optimal bit-complexity theorem.
\end{corollary}
\begin{proof}
By \eqref{eq:quadratic-margin} and $N-1\ge N/2$,
\[
 \Delta_N(a/N)\ge\frac{a(N-1)}{2(1+a)}\ge\frac{aN}{5}.
\]
The mean perturbation in \eqref{eq:normalized-feasibility} is at most $aN/64$, strictly less than this margin. The risk assertion follows from $J_N(a/N)\le a$ and Lemma~\ref{lem:resolvent-error}. Report averaging and rounding before regularization, as in Lemma~\ref{lem:finite-routing}, preserve the floor. For $1<N<2$ one may instead use \eqref{eq:exact-envelope-certificate}; no uniform small-margin assertion is being omitted at $N=1$.
\end{proof}

\section{Actual stopped mass and matching geometry}\label{sec:profile-proof}
The independent-report specialization below is the constrained-frontier mechanism already proved in~\cite{GTFStopped}; the comparison with correlated, controlled, model-indexed conditional rows is the extension. We give the argument in full because the history quantifiers and stopping-cost accounting are essential.
The reference measure in this section is a one-call acquired law, obtained from the reference preparation in Definition~\ref{def:recoverable}. It is not an assumption that any stopped policy has this law.

\begin{lemma}[Compact mass-constrained quantization]\label{lem:profile}
For every compact $K\subset\R^k$, probability $\nu$ supported on $K$, $m\ge1$ and real $t\ge1$, the minimum in \eqref{eq:profile} exists. For fixed centers put $d_C(x)=\min_j\norm{x-c_j}$. A minimizing density may be chosen in the form
\begin{equation}\label{eq:threshold-density}
 h(x)=t\one_{\{d_C(x)^2<\lambda\}}
       +t\xi\one_{\{d_C(x)^2=\lambda\}},\qquad 0\le\xi\le1,
\end{equation}
with $\int h\,d\nu=1$. The case $t=1$ permits $h=1$.
\end{lemma}
\begin{proof}
Centers may be projected onto the compact convex set $\conv K$, which cannot increase any distance to $K$. The set of densities $0\le h\le t$, $\int h=1$, is weak-star compact in $L^\infty(\nu)$. If centers converge, their squared nearest-center distance functions converge uniformly on $K$. The objective is therefore jointly continuous for this convergence and weak-star convergence of $h$. Compactness gives a minimum.

For fixed centers choose $\lambda$ such that
\[
 t\nu(d_C^2<\lambda)\le1\le t\nu(d_C^2\le\lambda)
\]
and choose $\xi$ to fill the possible boundary mass. Against any other admissible density $\widetilde h$, one has
\[
 \int(d_C^2-\lambda)(\widetilde h-h)\,d\nu\ge0,
\]
because on the sublevel set $h=t$ and on the strict superlevel set $h=0$. The equal total masses remove the $\lambda$ term. This proves optimality, including atoms and positive boundary mass. At a joint minimum the centers can also be replaced by the conditional means of their selected Voronoi cells, with any fixed Borel rule for ties; the squared-error projection identity proves that this does not increase the objective.
\end{proof}

\begin{lemma}[Acquisition mass bound]\label{lem:stopped-mass}
For any model-independent adaptive history policy and randomized stopping rule with $\E_\theta\tau\le N$, under the upper comparison in Definition~\ref{def:recoverable},
\begin{equation}\label{eq:actual-terminal-density}
 \law_\theta(X_\tau)\le b_*N\nu.
\end{equation}
More precisely, the unnormalized law of all acquired responses before and at stopping satisfies
\begin{equation}\label{eq:all-acquired}
 \E_\theta\sum_{t=1}^{\tau}\one_{\{X_t\in A\}}
       \le b_*\nu(A)\E_\theta\tau
       \quad(A\text{ Borel}).
\end{equation}
\end{lemma}
\begin{proof}
The event $\{\tau\ge t\}$ is determined before acquisition $t$, including all preceding private coins. Conditional on this information and the chosen action, the law of $X_t$ is at most $b_*\nu$. Summing
\[
 \E_\theta[\one_{\{\tau\ge t\}}\one_{\{X_t\in A\}}]
       \le b_*\nu(A)\Pp_\theta(\tau\ge t)
\]
over $t\ge1$ proves \eqref{eq:all-acquired} by Tonelli. The terminal indicator is no larger than the acquired-count sum. Its total mass is one, because finite mean implies almost-sure stopping, which proves \eqref{eq:actual-terminal-density}.
\end{proof}

\begin{lemma}[One-live-state realization]\label{lem:rejection}
Let $a_*N\ge1$. Given centers $C=\{c_1,\ldots,c_m\}$ and $0\le h\le a_*N$ with $\int h\,d\nu=1$, execute the fixed action $a_0$ and, after reading $x=f_{a_0}(y)$, accept with probability
\begin{equation}\label{eq:accept}
 A(x)=\frac{h(x)}{a_*N}.
\end{equation}
On acceptance, enter the nearest-center terminal label; otherwise return to the single live label. For every model this program satisfies
\begin{equation}\label{eq:rejection-bounds}
 \E_\theta\tau\le N,\qquad
 \E_\theta d_C(X_\tau)^2
       \le\frac{b_*}{a_*}\int d_C^2h\,d\nu.
\end{equation}
If all conditional response laws equal $\nu$, the choice $A=h/N$ with $0\le h\le N$ gives $\E\tau=N$ and $\law(X_\tau)=h\nu$ exactly.
\end{lemma}
\begin{proof}
At each surviving history the acceptance probability is at least
\[
 a_*\int\frac{h}{a_*N}\,d\nu=\frac1N.
\]
Consequently $\Pp_\theta(\tau>n)\le(1-1/N)^n$ and $\E_\theta\tau\le N$. The conditional expected squared distortion contributed by acceptance at that call is at most
\[
 \frac{b_*}{a_*N}\int d_C^2h\,d\nu.
\]
Summing it over surviving histories gives the second inequality of \eqref{eq:rejection-bounds}. This proof does not assume that repeated reports are independent or that the conditional law after rejection equals the original law.

For independent reference rows, the probability of acceptance at each surviving call is exactly $1/N$, and its accepted sublaw is $h\nu/N$. The sum of the geometric series is $N$, proving both exact assertions. Equivalently, the conditional response kernel being $\nu$ at every history suffices; no additional hidden-state independence is used.
\end{proof}

\begin{proof}[Proof of Theorem~\ref{thm:main}(ii)]
The protected audit identity gives, for every legal stopped policy and decoder,
\begin{equation}\label{eq:protected-pythagoras}
 \E_\theta\norm{Z-r_J}^2=B+\E_\theta\norm{X_\tau-r_J}^2.
\end{equation}
For a checkpoint there are at most $m$ decoder values. For an autonomous observer there are at most $M-1$, since at least one state is live. Conditional randomization of the decoder cannot improve squared loss; its mean belongs to $H$. Projecting centers onto $\conv K$ can only help. Lemma~\ref{lem:stopped-mass} bounds the actual terminal density by $b_*N$, so pointwise nearest-center comparison proves \eqref{eq:profile-check-lower}--\eqref{eq:profile-aut-lower} for every model and every randomized observer. It therefore proves the robust lower as well.

Apply Lemma~\ref{lem:rejection} to a minimizer in Lemma~\ref{lem:profile}, with $m=M-1$. This is a common program, determined by $\nu$, the comparison constants and its centers, not by $\theta$. It proves \eqref{eq:profile-aut-upper}. The same construction is available in the more permissive checkpoint architecture. If the conditional laws all equal $\nu$, the terminal law is exactly the minimizing $h\nu$, and the upper equals the lower. This proves \eqref{eq:exact-profile}. The stop is executed before the audit; no future audit information enters the accepting set.
\end{proof}

\begin{lemma}[Regular-support profile]\label{lem:regular}
Under \eqref{eq:regularity},
\begin{equation}\label{eq:profile-rate}
 c\,(tm)^{-2/s}\le\calQ_m^\nu(t)\le C\,(tm)^{-2/s}
 \quad(t\ge1,\ m\ge1),
\end{equation}
with constants depending only on the data specified in Theorem~\ref{thm:main}(iii).
\end{lemma}
\begin{proof}
For any center $c\in\conv K$, if $B(c,r)$ meets $K$, choose $x\in K\cap B(c,r)$. Then $K\cap B(c,r)\subset B(x,2r)$. Hence, when $2r\le r_0$, any union of $m$ such balls has $\nu$-mass at most $mC_0(2r)^s$. An admissible density is at most $t$. Choose a fixed sufficiently small $c_1>0$ and set $r=c_1(tm)^{-1/s}$, so that this acquired mass is at most $1/2$ and $2r\le r_0$. At least half the selected probability lies at distance at least $r$, proving the lower bound.

For the upper we record the covering and packing steps. A maximal $r$-separated subset of $K$ covers $K$ by $r$-balls, and its disjoint $r/3$-balls have mass at least $c_0(r/3)^s$. Thus its cardinality is at most $3^s/(c_0r^s)$. Enlarging constants to include the finitely small-cardinality regime gives a global cover by at most $m$ balls of radius $C_1m^{-1/s}$.

Conversely, choosing $\delta=c_2m^{-1/s}$ with $c_2$ sufficiently small, fewer than $m$ balls of radius $\delta$ cannot cover $K$, by the upper mass bound. A greedy construction therefore produces $m$ centers separated by at least $\delta$. For $t\ge t_0$, where $t_0$ depends only on the regularity constants, set
\[
 r=(c_0mt)^{-1/s}\le\delta/3.
\]
The $m$ disjoint $r$-balls have union $U$ of mass at least $mc_0r^s=1/t$. The density $h=\one_U/\nu(U)$ is admissible and its distance to these centers is at most $r$. This gives the required upper for $t\ge t_0$. For $1\le t<t_0$, use $h=1$ and the global $m$-center cover: $C_1^2m^{-2/s}\le C_1^2t_0^{2/s}(tm)^{-2/s}$. All changes for radii exceeding $r_0$ are absorbed by the fixed diameter constant, not by $t$ or $m$.
\end{proof}

\begin{proof}[Proof of Theorem~\ref{thm:main}(iii)]
For $a_*N\ge1$, insert Lemma~\ref{lem:regular} into the two profile bounds. For $1\le N<1/a_*$, stop after one report of $a_0$ and use the global $m$-center cover. This has mean acquisition one and excess at most $Cm^{-2/s}$, which is at most $Ca_*^{-2/s}(Nm)^{-2/s}$ in this range. The lower profile estimate holds for all $N\ge1$. The state deduction $m=M-1$ is exact throughout.
\end{proof}

\begin{proposition}[Atoms retain their physical mass]\label{prop:atoms}
For any compact acquired reference law and $m\ge1,t\ge1$,
\begin{equation}\label{eq:atom-threshold}
 \calQ_m^\nu(t)=0
 \quad\Longleftrightarrow\quad
 \text{some set of at most $m$ atoms has total $\nu$-mass at least $1/t$.}
\end{equation}
In particular, an atom of mass $w>0$ gives zero excess in the independent-reference experiment once $N\ge1/w$, with one live and one terminal state.
\end{proposition}
\begin{proof}
At a minimum with zero distortion, the probability measure $h\nu$ is supported on the at most $m$ centers. It can carry unit mass only if their $\nu$-mass is at least $1/t$. Conversely, restrict and normalize $\nu$ to such a finite set; the resulting density is at most $t$ and its distortion at those centers is zero. The causal assertion follows from the exact rejection implementation. All rejected observations are included in the mean $1/w$.
\end{proof}

\section{Two raw-kernel realizations}\label{sec:realizations}
The next results verify the preceding general theorem from physical kernels. Their finite affine representations are derived here and are not used as premises in Sections~\ref{sec:occupation-proof}--\ref{sec:profile-proof}.

\subsection{Correlated controlled geometric observations}
Fix $d\ge1$, let $K\subset[-1/2,1/2]^d$ be compact, and let $\nu$ be a centered probability with support $K$. The physical state space is the whole box; preparation at the origin is allowed even when the origin is outside $K$. A finite set of controls has unknown, fixed coefficients $\theta_a\in[-\eta,\eta]$, where $0\le\eta<1$. One raw acquisition from physical state $x$ selects the new state $y$ and reports it exactly, with kernel
\begin{equation}\label{eq:classical-kernel}
 K_a^\theta(dx',dy\mid x)
 =k_a^\theta(x,y)\,\delta_y(dx')\nu(dy),\qquad
 k_a^\theta(x,y)=1+\frac{4\theta_a}{d}\ip{x}{y}.
\end{equation}
Centering normalizes the kernel and $1-\eta\le k_a^\theta\le1+\eta$. The next report is generally correlated with the previous physical state. Different control coefficients change its conditional law. No reset to $\nu$ is made after a rejection. Preparation at $x=0$ produces $\nu$ in one actual call, uniformly in $\theta$.

The protected audit has conditionally independent coordinates $Z_i\in\{-1,0,1\}$, with
\begin{equation}\label{eq:classical-audit}
 \Pp(Z_i=\pm1\mid x)=\frac{x_i^2+1/2\pm x_i}{2},
 \qquad \Pp(Z_i=0\mid x)=\frac12-x_i^2.
\end{equation}
All these probabilities are positive on the box. Direct calculation gives $\E[Z_i\mid x]=x_i$ and $\operatorname{Var}(Z_i\mid x)=1/2$.

\begin{theorem}[Classical realization]\label{thm:classical}
The prepared kernels \eqref{eq:classical-kernel}--\eqref{eq:classical-audit} satisfy the finite affine and model-independent hypotheses of Theorem~\ref{thm:main}(i), and Definition~\ref{def:recoverable} with $a_*=1-\eta$, $b_*=1+\eta$ and $B=d/2$. If $\nu$ satisfies \eqref{eq:regularity}, then the same theorem yields, for the entire controlled unknown-coefficient family,
\begin{equation}\label{eq:classical-rate}
 \calR_\aut^\Theta(N,M)-\frac d2
       \asymp \bigl(N(M-1)\bigr)^{-2/s}.
\end{equation}
For uniform measure on the box, $s=d$. Every mean-budget lower holds for arbitrary history-dependent sensing and randomized stopping, not only for the implementing fixed action.
\end{theorem}
\begin{proof}
Let $\phi(x)=(x^\alpha)_{\alpha\in\{0,1,2\}^d}$, with the constant coordinate first, and let $D=\conv\phi([-1/2,1/2]^d)$. Work on its linear span with cone generated by $D$ and mass equal to the constant coordinate. The base is compact; the cone is closed and pointed. With $z_{e_i}$ denoting a first-moment coordinate, define the raw instrument
\begin{equation}\label{eq:classical-instrument}
 L_a^\theta(y)z=
   \left(z_0+\frac{4\theta_a}{d}\sum_i y_i z_{e_i}\right)\phi(y).
\end{equation}
It is positive on the cone because its scalar coefficient is the integral of the positive density \eqref{eq:classical-kernel} under a positive mixture of physical states. It is integrable and normalized. Its normalized output is $\phi(y)$, independently of the incoming predictive mixture. This verifies both finite affine update closure and report recoverability.

Every audit event has a probability polynomial of degree at most two in each coordinate, hence an affine evaluation on $D$. Conversely, the three univariate audit probabilities span $1,x_i,x_i^2$. Their products show that the finite audit-event family separates every coordinate of $D$. Future finite report-and-action words are obtained by positive compositions and integrations of \eqref{eq:classical-instrument}. Thus this is a finite closure of executable tests, not an arbitrary feature matrix. The task map extracts the first moments, and its second-moment evaluation is
\[
 c(\phi(x))=\frac d2+\sum_i x_i^2.
\]
On the actually acquired output $\phi(y)$, the conditional Bayes risk is $d/2$.

The raw likelihood bound $1-\eta\le k\le1+\eta$ survives conditioning on any history, since conditional mixtures preserve these measure inequalities. It holds for every adaptive choice of the next action. The first-call reference law is exactly $\nu$, and every fixed action is suitable for the upper construction. For completeness, the report kernel is continuous in the physical input, with
\[
 \norm{K_a^\theta(x,\cdot)-K_a^\theta(x',\cdot)}_\TV
 \le\frac{2\eta}{d}\int|\ip{x-x'}y|\,\nu(dy).
\]
On the full convex base the analogous bound uses the difference of first moments. Same-report predictive updates have zero discrepancy after normalization. These continuity and stability statements are consequences of the raw formula; neither a Fisher metric nor an assumed covariance geometry enters. Theorem~\ref{thm:main} now gives \eqref{eq:classical-rate}.
\end{proof}

\subsection{Noncommuting covariant quantum instruments}
Let $q\ge2$. A physical state is a density operator on $\mathbb C^q$. Let $\nu$ be the unitarily invariant probability on rank-one projections $P_\psi$. It is a probability on complex projective space, not on vectors with an observable phase. Each control selects a channel
\begin{equation}\label{eq:quantum-channel}
 \mathcal T_a^\theta(\rho)=\eta_a^\theta\frac{I}{q}\tr\rho
            +(1-\eta_a^\theta)U_a^\theta\rho(U_a^\theta)^*,
 \qquad \eta_a^\theta\in[\eta_0,1],\quad \eta_0>0,
\end{equation}
where the fixed coefficients and unitaries may be unknown. The raw instrument is
\begin{equation}\label{eq:quantum-instrument}
 L_a^\theta(d\psi)\rho
   =qP_\psi\mathcal T_a^\theta(\rho)P_\psi\,\nu(d\psi)
   =q\tr(P_\psi\mathcal T_a^\theta(\rho))P_\psi\,\nu(d\psi).
\end{equation}
It is a completely positive instrument: composition with $\mathcal T_a^\theta$ is followed by the positive measurement map $\rho\mapsto qP_\psi\rho P_\psi$. Since $\int P_\psi\,d\nu=I/q$, its integrated trace is preserved. The report is $P_\psi$ and the normalized output state is $P_\psi$.

After stopping, a protected covariant measurement on that output state reports $P_\zeta$ with density $q\tr(\rho P_\zeta)$ relative to $\nu$. The numerical audit vector, regarded in the real Hilbert space of Hermitian matrices, is
\begin{equation}\label{eq:quantum-audit}
 Z=(q+1)P_\zeta-I.
\end{equation}
The task is its Hilbert--Schmidt squared prediction loss. No further copy or quantum memory is available to the observer.

\begin{theorem}[Quantum realization]\label{thm:quantum}
The prepared instrument \eqref{eq:quantum-instrument} and audit \eqref{eq:quantum-audit} verify Theorem~\ref{thm:main}(i)--(iii), with
\[
 a_*=\eta_0,\qquad b_*=q,\qquad B=q^2+q-2,
 \qquad s=2(q-1).
\]
Preparation at $I/q$ realizes the reference law $\nu$ in one call for every model. Uniformly over the allowed channel family,
\begin{equation}\label{eq:quantum-rate}
 \calR_\aut^\Theta(N,M)-(q^2+q-2)
       \asymp \bigl(N(M-1)\bigr)^{-1/(q-1)}.
\end{equation}
The instrument has continuously many, generally noncommuting, state-dependent effects. No shared eigenbasis or finite designed measurement ensemble is required.
\end{theorem}
\begin{proof}
Use the real space of Hermitian matrices, its positive semidefinite cone and trace as mass. Positivity, compactness of the density-operator base and finite-dimensional affine closure follow directly from \eqref{eq:quantum-instrument}. The likelihood obeys
\[
 \eta_0\le q\tr(P_\psi\mathcal T_a^\theta(\rho))\le q
\]
for every density operator, hence for every incoming posterior mixture and adaptive history. The normalized map is the reported $P_\psi$, so it is report-recoverable. The action family changes probabilities, not the observer's access to the model.

We give the audit and geometry calculations. Unitary invariance, or the fourth moments of a normalized complex Gaussian vector, gives
\begin{equation}\label{eq:haar-second}
 \int P_\zeta\otimes P_\zeta\,\nu(d\zeta)
       =\frac{I+\mathsf{Swap}}{q(q+1)}.
\end{equation}
To verify this formula directly, independent standard complex Gaussian coordinates have independent exponential squared moduli. Their normalized squared moduli have uniform simplex density, with second moments $2/(q(q+1))$ on the diagonal and $1/(q(q+1))$ off the diagonal; independent phases annihilate the remaining entries. Pairing \eqref{eq:haar-second} with $\rho$ yields
\[
 \E[P_\zeta\mid\rho]=\frac{\rho+I}{q+1},\qquad \E[Z\mid\rho]=\rho.
\]
Moreover, $Z$ has eigenvalues $q,-1,\ldots,-1$, so $\norm Z_\HS^2=q^2+q-1$. On each acquired pure state, $\norm\rho_\HS^2=1$, giving constant Bayes risk $B=q^2+q-2$. The audit separates all density operators, since its mean is $\rho$. It therefore supplies executable informationally complete tests, rather than an assumed state metric.

For a fixed ray $P_\psi$, put $V=|\ip\psi\zeta|^2$. The same exponential-coordinate calculation gives
\[
 \Pp(V\ge u)=(1-u)^{q-1},\quad 0\le u\le1.
\]
Since $\norm{P_\psi-P_\zeta}_\HS^2=2(1-V)$, the actual reference measure has the exact ball mass
\begin{equation}\label{eq:projective-mass}
 \nu\{P_\zeta:\norm{P_\psi-P_\zeta}_\HS\le r\}
       =\left(\frac{r^2}{2}\right)^{q-1},\quad 0\le r\le\sqrt2.
\end{equation}
This proves the required regularity on the acquired projective support. In particular, its dimension is $2(q-1)$, not the ambient real dimension $q^2$ of Hermitian coordinates.

Finally, trace-norm contraction of a channel and the elementary bound $|\tr(PA)|\le\norm A_1$ give report-law continuity,
\[
 \norm{p_{a,\rho}^\theta-p_{a,\sigma}^\theta}_\TV
 \le\frac q2\norm{\rho-\sigma}_1.
\]
The posterior for a common report is exactly the same projection. Thus the raw likelihood, actual mass, global cover, update stability and protected audit conditions have all been verified. Apply Theorem~\ref{thm:main}.
\end{proof}

\begin{remark}[The role of physical preparation]
The quantum output projection and the classical next state are states of the experiment, not extra registers of the observer. The observer retains only its current classical label. A call to the quantum instrument includes the displayed controlled channel and measurement; the number of such calls is $\tau$. Gate counts, channel preparation, the protected audit call and the fixed apparatus dimension are additional physical costs. They are not equated with $\lceil\log_2M\rceil$ classical bits. The rate is a statement for this declared instrument interface, not for arbitrary laboratory implementations with free resets.
\end{remark}

\section{Informative instruments without posterior contraction}\label{sec:effective-realizations}
The preceding classical and quantum realizations yield explicit geometric rates because the current task response is report-recoverable. The exact-budget construction has a different, strictly broader hypothesis. We now verify it on informative experiments in which the current report does not determine the posterior. Neither proof below invokes a realization-specific theorem to establish the general construction.

\subsection{Controlled Gaussian hidden-state experiments}
Let the hidden state set be $\{1,\ldots,d\}$ and let the action set be finite. An action first makes a transition $P_a^\theta(j\mid i)$ and then reports a vector $y\in\R^p$ with density $g_{aj}^\theta(y)$ of $\mathcal N(m_{aj}^\theta,\Sigma_{aj}^\theta)$. Thus the raw joint kernel and its unnormalized instrument are
\begin{equation}\label{eq:effective-gaussian}
 K_a^\theta(j,dy\mid i)=P_a^\theta(j\mid i)g_{aj}^\theta(y)\,dy,
 \qquad L_a^\theta(y)_{ji}
       =P_a^\theta(j\mid i)g_{aj}^\theta(y)/w(y),
\end{equation}
where $\lambda(dy)=w(y)dy$ is a fixed nondegenerate Gaussian probability with explicitly known density, chosen independently of $\theta$. The ordered physical space is $\R^d$ with the nonnegative cone and sum as mass. A prepared hidden state has distribution $x_0^\theta$. After stopping, an independent protected audit returns the one-hot vector of the current hidden state. Its mean is the posterior vector and its second moment is one. Decoders lie in the simplex, and the squared loss is bounded by $2$.

\begin{theorem}[Exact-budget finite construction for informative Gaussian kernels]\label{thm:effective-gaussian}
Suppose the family of $x_0^\theta,P_a^\theta,m_{aj}^\theta,\Sigma_{aj}^\theta$ is effectively compact and computable, with
\[
 \norm{m_{aj}^\theta}\le C_m,\qquad
 c_\Sigma I\le\Sigma_{aj}^\theta\le C_\Sigma I,
 \qquad 0<c_\Sigma\le C_\Sigma<\infty.
\]
Assume the supplied report digitization can implement bounded-box partitions with a certified vanishing boundary error, and records the entire complement as a charged tail outcome. Then Theorems~\ref{thm:main}(i) and \ref{thm:effective-budget} hold for every rational $N\ge1$ and $M\ge2$.
No positivity floor, invertibility, irreducibility, or mixing rate is required for $P_a^\theta$. The theorem covers identity and reducible transitions and unknown fixed parameters using one common program.
\end{theorem}
\begin{proof}
Each matrix in \eqref{eq:effective-gaussian} is positive, its integrated column sums are one, and its integral operator norm is finite. The protected one-hot audit gives every coordinate of the posterior as an executable future-test evaluation. These evaluations separate the simplex. All unnormalized report pullbacks remain in the $d$-dimensional linear space. This verifies the finite predictive representation from the raw experiment. If a smaller test family is used, its invariant separating span, rather than the whole simplex, is the relevant quotient. Here the audit already separates all coordinates.

We verify the uniform approximation certificate. For $R\ge2C_m+1$, Gaussian tails and the covariance bounds give computable constants $C,c>0$ such that
\begin{equation}\label{eq:gaussian-tail-compile}
 \sup_{\theta,a,j}\int_{\norm y_\infty>R}g_{aj}^\theta(y)dy
        \le C\exp(-cR^2).
\end{equation}
For example, apply the one-dimensional Gaussian Chernoff bound to each coordinate and take the union bound. Partition the box $[-R,R]^p$ into half-open rational cubes of side at most $h$, and keep its complement as one additional cell. On the box, $w$ has a positive computable minimum. Differentiating the Gaussian formula shows that $g_{aj}^\theta/w$ has a computable uniform gradient bound $C_R$. The mean-value theorem bounds the oscillation of every instrument entry on a cube by $C_R\sqrt p\,h$. In a fixed finite dimension, summing entry bounds bounds the induced $\ell^1$ operator norm, so the integrated interior approximation error is at most $C'_Rh$.

On the tail cell, positivity gives
\[
 \int_{\mathrm{tail}}\bnorm{L-\overline L}\,d\lambda
    \le2\int_{\mathrm{tail}}\bnorm L\,d\lambda
    \le C'\exp(-cR^2).
\]
For the last inequality, the maximum column sum is at most the sum of all entries, bounded by a fixed-dimensional sum of Gaussian tails. Consequently
\begin{equation}\label{eq:gaussian-partition-modulus}
 \eta_{R,h}\le C'_Rh+C'\exp(-cR^2).
\end{equation}
First choose $R$ and then $h$. The original tail remains part of every physical call and of every cell integral; it is neither discarded nor conditioned away. A report-digitization error is added to this bound as specified in Definition~\ref{def:effective}.

For model continuity, derivatives of a Gaussian density with respect to its mean and covariance are that density times a polynomial of degree at most two in $y-m$, with coefficients bounded by powers of $c_\Sigma^{-1}$. Their absolute integrals are uniformly bounded on the stated parameter body. Integrate these derivatives along a straight segment in mean--covariance space; the covariance segment retains the eigenvalue bounds. This proves a computable $L^1$ modulus for the densities. Matrix multiplication by stochastic $P_a^\theta$ and finite-dimensional triangle inequalities give the required integral instrument modulus, even when some transition entries vanish. The preparation modulus is part of the effectively compact data.

Box-cell integrals are computable uniformly: truncate an already certified tail and integrate the smooth bounded density on a compact box with its explicit derivative bounds. The complementary cell equals the normalization minus the interior cells. Simplex decoder nets are finite and rational. Matrix products and bounded linear audit costs have certified interval evaluations. This proves all of Definition~\ref{def:effective}. The general theorems apply. The proof never bounds a normalized likelihood ratio by a positive constant, and never reconstructs an exact posterior in the observer.
\end{proof}

The identity-transition member is a standard static hidden-state Gaussian sensing experiment, and positive-transition members are informative nonlinear filters. The theorem provides a two-sided finite certificate for their hard-state mean-budget optimum; it does not claim an explicit sharp exponent for all of them, nor an efficient solution of unknown-kernel learning. The parameter family is known to the program designer; the true parameter is not.

\subsection{Noncommuting quantum instruments, including singular effects}
Let $E$ be the real Hermitian matrices on $\mathbb C^q$, with the positive semidefinite cone, trace mass, and trace base norm. Let $Y$ be a compact metric report space with a fixed computable probability measure $\lambda$ and computable finite partitions of vanishing diameter. For each action and parameter, let $E_a^\theta(y)$ be positive semidefinite effects satisfying
\begin{equation}\label{eq:effective-quantum-effects}
 \int E_a^\theta(y)\lambda(dy)=I,\qquad
 0\le E_a^\theta(y)\le C_EI.
\end{equation}
Let $\mathcal T_a^\theta$ be a completely positive trace-preserving channel. The physical instrument is
\begin{equation}\label{eq:effective-quantum}
 L_a^\theta(y)\rho=
 \mathcal T_a^\theta\bigl(E_a^\theta(y)^{1/2}\rho E_a^\theta(y)^{1/2}\bigr).
\end{equation}
Take an effectively compact, uniformly computable family with computable uniform moduli for effects in $(\theta,y)$ and channels in the induced trace norm. Effects may lose rank. No simultaneous diagonalization is assumed. A protected informationally complete finite audit is fixed once and for all: a finite informationally complete POVM $\{F_b\}$ with real dual matrices $D_b$ satisfies $\sum_b\tr(F_b\rho)D_b=\rho$, and reports $D_b$. It has bounded, affine first and second moments. A finite such POVM can be obtained from sufficiently small positive perturbations of $I$ in a Hermitian basis and their negative partners; its finite spanning measurement matrix has a computable right inverse. Decoders lie in the density-operator body.

\begin{theorem}[Exact-budget quantum construction at changing rank]\label{thm:effective-quantum}
Under these raw hypotheses, including certified finite report partition execution, Theorem~\ref{thm:effective-budget} gives two-sided finite certificates for every rational $N\ge1$ and $M\ge2$. The observer is classical, its stopping and decoder labels are charged in $M$, and it uses no retained quantum memory. Rank-deficient preparations, rank-changing effects, and noncontracting unitary channels are allowed.
\end{theorem}
\begin{proof}
Complete positivity of \eqref{eq:effective-quantum} and integrated trace normalization follow from \eqref{eq:effective-quantum-effects}. The audit mean is $\rho$ and separates all density operators. Its finite set of first and second moments belongs to the Hermitian linear space; report pullbacks are linear in that same space. Thus the finite executable-test representation is verified without assuming a Fisher matrix or a smooth state parametrization.

For positive matrices $A,B$, the operator-norm square-root inequality
\begin{equation}\label{eq:sqrt-holder}
 \norm{A^{1/2}-B^{1/2}}_\infty\le\norm{A-B}_\infty^{1/2}
\end{equation}
is valid including at zero eigenvalues. One elementary proof puts $t=\norm{A-B}_\infty$ and uses $A\le B+tI$, operator monotonicity of the square root, and $(B+tI)^{1/2}\le B^{1/2}+\sqrt t I$; exchange $A,B$ for the opposite order. Operator monotonicity follows, for instance, by integrating $x/(x+s)$ against its positive square-root integral representation. These are finite-matrix statements.

For Hermitian $X$ with trace norm at most one, the ideal property of the trace norm and a two-term expansion give
\[
 \norm{A^{1/2}XA^{1/2}-B^{1/2}XB^{1/2}}_1
 \le(\norm{A^{1/2}}_\infty+\norm{B^{1/2}}_\infty)
                  \norm{A^{1/2}-B^{1/2}}_\infty.
\]
Trace-preserving positive maps contract trace norm on Hermitian inputs. Consequently effect and channel discrepancies $u,v$ give a raw instrument discrepancy at most
\begin{equation}\label{eq:quantum-raw-modulus}
 2\sqrt{C_E}\sqrt u+C_Ev.
\end{equation}
The same estimate in the report coordinate supplies a uniform oscillation bound on every partition cell, hence \eqref{eq:report-approximation}. No inverse of a density operator, normalized posterior, or effect is used. Continuity is therefore retained across rank changes.

Integrals of the computable uniformly continuous finite-dimensional matrix functions are computable with certified partition error under the stipulated computable measure. Density-operator decoder nets can be constructed by approximating a square-root factor with rational real and imaginary entries, forming $BB^*/\tr(BB^*)$, and taking sufficiently fine bounded factor grids; these are legal, finitely described positive matrices and cover the compact state body. The finite audit moments and all product evaluations are computable. Definition~\ref{def:effective} follows, so the general theorem proves the assertion.
\end{proof}

A concrete continuously indexed family takes $Y$ to be a compact symmetric trace-zero Hermitian ball, $E^\theta(H)=I+\eta_\theta H$ with $0\le\eta_\theta\le1$ and $\norm H_\infty\le1$, and any effectively compact channel family. When $\eta_\theta=1$ and $H$ has eigenvalue $-1$, an effect is singular. A symmetric computable measure on $Y$ has mean zero, so normalization holds. Discrete, Cantor-parametrized, or finite mixtures of such report measures also qualify when the specified partitions and integrals are computable. The actual weights remain in $\lambda$ and in every instrument integral; they are not replaced by ambient volume. Report support may vary with the incoming state through zero effects. A single dominating reference measure remains an assumption.

\subsection{The same normalized transfer in both realizations}
\begin{corollary}[Raw verification of the sharp-order transfer]\label{cor:raw-normalized}
Consider two Gaussian experiment families with the common compact parameter and report interfaces of Theorem~\ref{thm:effective-gaussian}. Suppose their preparations, transition probabilities, Gaussian means and covariance matrices differ by at most $u$, uniformly in the unknown model and action; retain the stated common ellipticity bounds. Then, for some constant $C$ depending only on these raw bounds, the same-budget risk difference is bounded by
\begin{equation}\label{eq:gaussian-normalized}
 \min\{L,\ L[2C(N+1)u+\sqrt{C^2(N+1)^2u^2+2c_NC(N+1)u}]+\alpha\}.
\end{equation}
Here $N\ge1,M\ge2$, and $\alpha$ denotes any additional uniform audit error.

For the quantum families of Theorem~\ref{thm:effective-quantum}, suppose preparations differ in trace norm by at most $u$, channels in the declared induced trace norm by at most $u$, and effects uniformly in operator norm by at most $u$. Put
\[
 z_u=u+CN(\sqrt u+u).
\]
The risk difference is at most
\begin{equation}\label{eq:quantum-normalized}
 \min\{L,\ L[2z_u+\sqrt{z_u^2+2c_Nz_u}]+\alpha\}.
\end{equation}
This conclusion permits rank changes. If both effect families instead have a common strictly positive lower bound, one may replace $\sqrt u+u$ by a constant multiple of $u$. These are sufficient moduli; the fractional effect-parameter exponent at rank change is not asserted sharp. Both realizations optimize over the entire declared common-program class at the same $N,M$.
\end{corollary}
\begin{proof}
For Gaussian densities, integrate the derivatives in the means and covariance entries along the line joining the two parameter values. Common ellipticity and bounded means bound these derivatives in $L^1$ by a constant. Multiplying by the transition probabilities and summing the finite input/output indices yields the integral induced base-norm instrument error $\eta\le Cu$. The preparation contributes at most $Cu$. The proof integrates over the whole report space; the tail is not discarded. Positivity implies the same bound after every common report-routing kernel. Apply Theorem~\ref{thm:normalized-transfer} with $\zeta\le C(N+1)u$.

In the quantum case, write $A=\sqrt E$, $B=\sqrt{\widetilde E}$. On a unit trace positive input $x$,
\[
 \|AxA-BxB\|_1\le(\|A\|+\|B\|)\|A-B\|.
\]
Use \eqref{eq:sqrt-holder}, the uniform effect bound and the probability reference measure, then add the channel error. Positive/negative decomposition extends the bound to the induced base norm, giving $\eta\le C(\sqrt u+u)$. This is an unnormalized-instrument estimate, not continuity of a posterior divided by a small likelihood. If the effects are bounded below by $\kappa I$, the identity
\[
 A(A-B)+(A-B)B=E-\widetilde E
\]
has the solution $A-B=\int_0^\infty e^{-tA}(E-\widetilde E)e^{-tB}\,dt$, whose operator norm is at most $u/(2\sqrt\kappa)$. This gives the improved linear bound. Theorem~\ref{thm:normalized-transfer} proves the assertions, including unchanged state and mean budgets and the separate audit error.
\end{proof}

\section{Singular acquisition and changing task strata}\label{sec:singular}
The singular independent-preparation and nonconstant-Bayes-risk mechanisms are inherited from~\cite{GTFStopped}. Here the regularity proof applies to all real mean budgets and the comparison survives correlated model families; the raw changing-rank instrument below supplies a separate realization of the weighted profile.
\subsection{A noninteger-dimensional raw realization}
Let $\mu_C$ be the centered uniform middle-thirds Cantor probability on $[-1/2,1/2]$, and let $\nu=\mu_C^{\otimes d}$. This is an actual preparation law for \eqref{eq:classical-kernel}; it is not Lebesgue measure restricted to a null set.

\begin{proposition}[Singular controlled acquisition]\label{prop:cantor}
For the classical raw experiment with $\nu=\mu_C^{\otimes d}$, Theorem~\ref{thm:main} gives
\begin{equation}\label{eq:cantor-rate}
 \calR_\aut^\Theta(N,M)-\frac d2
 \asymp \bigl(N(M-1)\bigr)^{-2/(d\log2/\log3)}.
\end{equation}
The conclusion holds for the full controlled correlated coefficient family, with all failed and accepted acquisitions charged.
\end{proposition}
\begin{proof}
A level-$n$ Cantor cylinder has length $3^{-n}$ and mass $2^{-n}$. If $3^{-(n+1)}<r\le3^{-n}$, a ball centered on the Cantor set contains the level-$(n+1)$ cylinder containing its center, so its mass is at least $2^{-(n+1)}$. An interval of radius $r$ meets at most a fixed number of level-$n$ cylinders, since these intervals have length $3^{-n}$ and disjoint interiors separated by gaps at least $3^{-n}$. Its mass is at most $C2^{-n}$. Thus $\mu_C(B(x,r))\asymp r^{\log2/\log3}$ uniformly in $x$ and $r$.

In $d$ dimensions, a Euclidean ball contains the product of coordinate balls of radius $r/\sqrt d$ and is contained in the product of coordinate balls of radius $r$. Multiplying the one-dimensional estimates proves \eqref{eq:regularity} with $s=d\log2/\log3$. The kernel, finite test closure, audit and model-uniform comparison were proved for any centered $\nu$ in Theorem~\ref{thm:classical}; no density relative to ambient volume was used there. The general regular-support theorem therefore proves \eqref{eq:cantor-rate}.
\end{proof}

A mixture $\nu=w\delta_0+(1-w)\mu_C^{\otimes d}$ is also centered and is an actual reference preparation in the same experiment. The mass-constrained profile, not a single regular exponent, applies. In the independent reference member, $\calQ_1^\nu(N)=0$ once $N\ge1/w$. The mean cost of selecting the atom is precisely $1/w$; its small physical probability is not erased by normalization. For the whole correlated family, the generic comparison gives a valid sufficient budget $1/(a_*w)$. In this particular centered kernel, the atom's next-call probability is exactly $w$ at every physical state, so direct rejection gives the sharper budget $1/w$ as well.

\subsection{Nonconstant conditional Bayes risk}
Changing-rank quantum outputs naturally have different irreducible variances. Subtracting one common number would change the problem. There is an exact extension that retains them. Let $X=f_a(Y)\in K$, and suppose the protected squared audit has conditional Bayes risk $\beta(X)$, where $\beta\ge0$ is known and continuous on compact $K$. Define
\begin{equation}\label{eq:weighted-profile}
 \calQ_{m,\beta}^\nu(t)=\min_{C,\,0\le h\le t,\,\int h=1}
       \int\bigl(\beta(x)+d_C(x)^2\bigr)h(x)\nu(dx).
\end{equation}

\begin{corollary}[Stratum-sensitive total risk]\label{cor:weighted-profile}
Under the report comparison hypotheses, with this nonconstant conditional Bayes risk, the checkpoint and autonomous \emph{total} risks obey the bounds of Theorem~\ref{thm:main}(ii), with $\calQ_m^\nu$ replaced by $\calQ_{m,\beta}^\nu$ and with no subtraction of $B$. For independent reference reports,
\begin{equation}\label{eq:weighted-exact}
 \calR_\cp(N,m)=\calQ_{m,\beta}^\nu(N),\qquad
 \calR_\aut(N,M)=\calQ_{M-1,\beta}^\nu(N).
\end{equation}
The model-uniform upper factor remains $b_*/a_*$ when $a_*N\ge1$.
\end{corollary}
\begin{proof}
Replace the nearest-center cost $d_C^2$ in Lemmas~\ref{lem:profile} and \ref{lem:rejection} by the nonnegative continuous cost $\beta+d_C^2$. Compactness and threshold selection are unchanged. Conditional audit projection gives total risk $\E[\beta(X_\tau)+\norm{X_\tau-r_J}^2]$. Lemma~\ref{lem:stopped-mass} proves the lower, and the same acceptance rule proves the upper and the independent-report equality. No comparison of different Bayes baselines has been made.
\end{proof}

\subsection{Raw quantum instruments with changing rank}
For $1\le r\le q$, let $\nu_r$ be Haar probability on the orbit
\[
 K_r=\{UP_rU^*/r:U\in\mathrm U(q)\},
\]
where $P_r$ is a fixed rank-$r$ projection. The case $r=q$ is the single state $I/q$. Choose fixed weights $w_r\ge0$, $\sum_rw_r=1$, with at least two positive weights. A raw call reports the marked state $(r,\sigma)$ and prepares $\sigma$ through
\begin{equation}\label{eq:rank-instrument}
 L_a^\theta(r,d\sigma)\rho
       =w_rq\tr(\sigma\mathcal T_a^\theta(\rho))\sigma\,\nu_r(d\sigma).
\end{equation}
This is a positive, indeed completely positive, measure-and-prepare instrument. Each effect is positive; the prepared state is positive; and $\int\sigma\,d\nu_r=I/q$ proves integrated trace normalization. The same protected covariant audit \eqref{eq:quantum-audit} has conditional mean $\sigma$ and
\begin{equation}\label{eq:rank-beta}
 \beta(\sigma)=q^2+q-1-\tr\sigma^2
              =q^2+q-1-\frac1r.
\end{equation}
The finitely many rank orbits are disjoint compact sets, and this function is continuous on their union. No observer register containing $r$ is required after the terminal label has been selected; $r$ is part of the current report only.

\begin{proposition}[Changing-rank acquired profile]\label{prop:rank}
The raw instrument \eqref{eq:rank-instrument} satisfies the general finite-affine theorem and Corollary~\ref{cor:weighted-profile} with actual reference law $\nu=\sum_rw_r\nu_r$, $a_*=\eta_0$ and $b_*=q$. In the independent member $\eta_a^\theta=1$ for all actions, its exact autonomous risk is
\begin{equation}\label{eq:rank-exact}
 \calR_\aut(N,M)
  =\min_{C,\,h}
       \sum_rw_r\int_{K_r}
        \left(q^2+q-1-\frac1r+d_C(\sigma)^2\right)
                          h(\sigma)\nu_r(d\sigma),
\end{equation}
where $0\le h\le N$ and $\sum_rw_r\int h\,d\nu_r=1$, and $|C|\le M-1$. This covers mixtures of continuous rank strata and a genuine full-rank atom without conditioning away any weight.
\end{proposition}
\begin{proof}
Use the positive semidefinite cone and trace. Formula \eqref{eq:rank-instrument} is linear in the incoming density matrix and has normalized output equal to the report. The density relative to $\nu$ lies between $\eta_0$ and $q$, because $\tr\sigma=1$ and $\mathcal T_a^\theta(\rho)\ge\eta_0I/q$. Preparation at $I/q$ gives exactly $\nu$. The audit calculations establishing its affine moments and separation are those already proved in Theorem~\ref{thm:quantum}, valid for every density operator. Report-kernel continuity follows by the same trace-norm bound, and same-report normalization is independent of the incoming state. Thus all hypotheses are verified from the raw maps. Corollary~\ref{cor:weighted-profile} proves \eqref{eq:rank-exact}.
\end{proof}

The formula optimizes the genuine task, including its different rank-dependent Bayes risks. The zero-distortion atom criterion in Proposition~\ref{prop:atoms} concerns the distortion component; it is not a claim that selecting the full-rank atom minimizes \eqref{eq:rank-exact}. This distinction is material when the stopping decision changes the irreducible audit risk.

\section{A sharp singular acquisition and memory law}\label{sec:singular-memory}
An acquired support may change when a positive but arbitrarily small preparation mass becomes visible. The mean budget then charges a long occupation on that component, even though a bounded terminal metric sees only its small probability. We give an exact finite-state law for this effect. It is a sharpness consequence of the general theory, not a replacement for the experiment-level theorem.

\subsection{The raw experiment and its actual occupations}
Fix $p\in(0,1)$ and write $q=1-p$. The physical states are $A$ (ready), $S$ (success) and $T$ (trapped). There is one sampling action and reports $f,s$. From $A$ the kernel sends the state to $S$ with report $s$ with probability $p$, and leaves it at $A$ with report $f$ with probability $q$. States $S,T$ are absorbing, with reports $s,f$, respectively. The preparation is
\begin{equation}\label{eq:rare-preparation}
 x_0^b=(1-b)e_A+be_T,\qquad 0\le b\le1.
\end{equation}
The raw instruments are nonnegative matrices on $\R^3$ with mass functional $\one^\top$. Positivity here means preservation of the positive cone, not a strictly positive lower bound on every transition. The next success report distinguishes $A$ from $T$, and a protected current audit distinguishes $S$ from them. These tests and the constant test span the affine prediction coordinates. Their raw positive update closes in three coordinates; a Fisher or covariance geometry is not an input.

We use two protected terminal tasks, scored only after the machine has committed to a terminal label. The first has the loss $\one_{\{X_\tau\ne S\}}$ and a singleton decoder body. The second audits a binary variable $Y$ that equals one at $S$ and is a fresh fair bit at $A,T$; the forecast lies in $[0,1]$ and the loss is $(Y-c)^2$. These are two separately declared tasks on the same sampling experiment. The audit costs one additional physical call and is not an observation available to the stopping rule.

For the failure task, stopping immediately on any success report never increases loss and can only decrease the mean call count. Merge all terminal labels. A machine using at most $M$ total states then has at most $k=M-1$ live states. Under a sequence consisting entirely of failure reports, its routing generates an absorption count $T_0\ge1$. If $b>0$, this is an \emph{actually prepared} component and $\E T_0<\infty$ follows from the mean budget. On the ready component the first success time $G_p$ is independent geometric with parameter $p$, independent of the routing coins. Hence, with
\[
 m=\E T_0,\qquad \varphi=\E q^{T_0},
\]
the actual failure probability and mean acquisition are exactly
\begin{equation}\label{eq:rare-flow}
 F=b+(1-b)\varphi,\qquad
 \E_b\tau=\frac{1-F}{p}+bm.
\end{equation}
The quantities $bm$ and $(1-F)/p$ are the actual trapped and ready live occupations. They are not normalized conditional substitutes for the preparation law.

\begin{theorem}[Exact memory frontier at every mean budget]\label{thm:rare-memory}
Let $b\in(0,1)$, $N\ge1$ and $M\ge2$, and put $k=M-1$. Let $g_{k,p}$ be the explicit discrete clock function in \eqref{eq:clock-frontier}. There is a unique $m_*\ge1$ satisfying
\begin{equation}\label{eq:rare-root}
 \frac{(1-b)(1-g_{k,p}(m_*))}{p}+bm_*=N.
\end{equation}
Define
\begin{equation}\label{eq:rare-profile}
 F_{p,b}(N,M)=b+(1-b)g_{k,p}(m_*).
\end{equation}
The optimal autonomous failure risk, over all homogeneous Borel randomized programs using at most $M$ total states, is exactly $F_{p,b}(N,M)$. It is attained by a clock with at most $k$ live states and one terminal state, with actual mean exactly $N$.

For the binary squared-prediction task on the same raw kernel,
\begin{equation}\label{eq:rare-prediction}
 \frac14 F_{p,b}(N,M)
 \le\calR_{\aut,\mathrm{pred}}^b(N,M)
 \le\frac12 F_{p,b}(N,M)-\frac14F_{p,b}(N,M)^2.
\end{equation}
The upper is implementable using the same clock and a single forecast. When $M=2$ it is the exact prediction risk. All statements hold also in the robust form with an unknown fixed preparation parameter in $[0,b]$, using one program designed for the known endpoint $b$.
\end{theorem}
\begin{proof}
The left side of \eqref{eq:rare-root}, regarded as a function of $m$, is continuous and strictly increasing from one to infinity, since $g_{k,p}$ is strictly decreasing and $b>0$. This proves existence and uniqueness.

Consider any competitor for the failure task. After the harmless modifications preceding the theorem, its all-failure clock has at most $k$ phases. Lemma~\ref{lem:discrete-clock} and \eqref{eq:rare-flow} imply
\[
 F\ge b+(1-b)g_{k,p}(m),\qquad
 m\le\frac{pN-1+F}{pb}.
\]
Since $g_{k,p}$ is decreasing,
\begin{equation}\label{eq:rare-fixed-inequality}
 F\ge b+(1-b)g_{k,p}\!\left(\frac{pN-1+F}{pb}\right).
\end{equation}
The argument is at least one for a feasible competitor. The right side is strictly decreasing in $F$, and equality occurs at \eqref{eq:rare-profile}. Thus \eqref{eq:rare-fixed-inequality} gives $F\ge F_{p,b}(N,M)$ for every competitor. The clock attaining $g_{k,p}(m_*)$ realizes equality in both equations of \eqref{eq:rare-flow}, proving the exact value and its state accounting.

For the prediction task, condition on the true terminal physical state as well as the terminal label. The fresh audit variance is $1/4$ on failure and zero on success. Every forecast, including a randomized one, therefore has risk at least $F/4$. Its failure probability cannot be smaller than the failure optimum already proved. For the upper use one terminal forecast $c=1-F_{p,b}(N,M)/2$. The audit has unconditional success probability $1-F_{p,b}(N,M)/2$, so its variance is the upper side of \eqref{eq:rare-prediction}. If $M=2$, the terminal forecast is necessarily a single constant (randomization cannot improve squared loss). Its best risk for failure probability $F$ is $F/2-F^2/4$, an increasing function on $[0,1]$. Stopping immediately on success leaves $F$ unchanged and reduces the mean. The failure optimum consequently also minimizes the prediction risk.

For the robust assertion, a fixed all-failure clock has $F_\theta=\theta+(1-\theta)\varphi$, increasing in $\theta$. Its mean is
$(1-\theta)\E\min(T_0,G_p)+\theta\E T_0$, also increasing in $\theta$. The endpoint construction is therefore feasible throughout $[0,b]$. Its chosen forecast is at least $1/2$, and its squared risk is increasing in the failure probability. The endpoint gives the lower bound on any common program. These facts prove the robust exact failure assertion, the prediction comparison, and the exact $M=2$ assertion without access to the true preparation parameter.
\end{proof}

At $b=0$ the failure optimum is $\max\{1-pN,0\}$. For $1\le N<1/p$ a one-phase randomized failure stop attains it; for $N\ge1/p$ waiting for the first success does so. Under this nominal preparation the all-failure path has no positive prepared mass, and its counterfactual clock may never halt. This does not contradict the preceding proof, which uses $b>0$ to charge its mean.

\begin{corollary}[Memory-dependent singular exponent]\label{cor:memory-exponent}
At the critical mean budget $N=1/p$, fix $k=M-1$. If
\begin{equation}\label{eq:phase-range}
 b\bigl(kp-1+q^k\bigr)\le q^k,
\end{equation}
then $m_*\ge k$ and the exact failure risk is the unique $F\in[b,1]$ solving
\begin{equation}\label{eq:memory-polynomial}
 F=b+(1-b)\left(\frac{qkb}{F+qkb}\right)^k.
\end{equation}
For each fixed $p\in(0,1)$ and $M\ge2$, as $b\downarrow0$,
\begin{equation}\label{eq:memory-exponent}
 F_{p,b}(1/p,M)\sim(qk)^{k/(k+1)}b^{k/(k+1)}.
\end{equation}
Consequently the prediction risk is of this same order, with constants between $1/4$ and $1/2$ times the failure scale. The exponent depends on paid finite-state memory although the physical carrier always has three coordinates. The asymptotic equivalence is for fixed $p,M$, not uniformly over a growing number of states.
\end{corollary}
\begin{proof}
The mean of the deterministic $k$-call failure clock is
$(1-b)(1-q^k)/p+bk$. Comparing it with $1/p$ yields exactly \eqref{eq:phase-range}. Thus the second branch of \eqref{eq:clock-frontier} applies. The mean identity in \eqref{eq:rare-flow} gives $m_*=F/(pb)$, and substitution proves \eqref{eq:memory-polynomial}.

The root tends to zero: for each fixed $f>0$ the right side of \eqref{eq:memory-polynomial} tends to zero. On the other hand $F/b\to\infty$, since a bounded subsequence of $F/b$ would leave its right side bounded away from zero. Rewrite the equation as
\[
 (F-b)(F+qkb)^k=(1-b)(qkb)^k.
\]
Dividing by $F^{k+1}$ and using $b/F\to0$ proves
\eqref{eq:memory-exponent}. The prediction comparison follows from
\eqref{eq:rare-prediction}.
\end{proof}

\begin{theorem}[Sharp preparation perturbation with two total states]\label{thm:rare-preparation}
For $M=2$ and $N=1/p$, the exact failure risk is
\begin{equation}\label{eq:two-state-preparation}
 F=\frac{pb+\sqrt{p^2b^2+4qb}}2,
 \qquad \rho_* =\frac{pb}{F}.
\end{equation}
An optimal program halts on success, and halts after each failure with probability $\rho_*$. Its exact prediction risk is $F/2-F^2/4$. The nominal value at $b=0$ is zero. The preparation base-norm distance is $\delta_0=2b$, while the instruments and task are unchanged. Thus the value difference is of order $\sqrt{c_N\delta_0}$ as $b\downarrow0$, for fixed $N>1$. For $N\ge2$ and $0<b\le1/4$, the comparison constants can be chosen numerically, independently of $N$.
\end{theorem}
\begin{proof}
For $k=1$, \eqref{eq:memory-polynomial} reduces to
$F^2-pbF-qb=0$. The positive root is \eqref{eq:two-state-preparation}. The trapped occupation is $b/\rho_*$ and, at mean equality, equals $F/p$, proving the routing formula. Theorem~\ref{thm:rare-memory} proves its global optimality and the prediction identity, including the robust interval assertion. The two preparations differ by $b(e_T-e_A)$, of base norm $2b$. Finally $\sqrt{qb}\le F\le pb+\sqrt{qb}$ and $F/4\le F/2-F^2/4\le F/2$. For $N\ge2$, $q\ge1/2$, giving the stated uniform comparisons.
\end{proof}

\subsection{Sharp one-call perturbation, including acquisition scaling}
Keep the preparation at $A$. Alter only its failure transition: from $A$, report $f$ and stay at $A$ with probability $q-h$, report $f$ and enter $T$ with probability $h$, and report $s$ and enter $S$ with probability $p$. Here $0<h<q$; the other transitions are unchanged. The joint raw-instrument base-norm distance from $h=0$ is exactly $\eta=2h$. In particular this is a primitive error, not a difference between normalized posterior parametrizations.

\begin{theorem}[Sharp one-call sensitivity at the active budget]\label{thm:rare-instrument}
At $N=1/p$ and $M=2$, put $a=(N-1)h$. The optimal failure risk over all legal randomized programs is
\begin{equation}\label{eq:two-state-instrument}
 F=\frac{\sqrt{a^2+4a}-a}{2},\qquad
 \rho_* =\frac{h(1-F)}{F+h(1-F)}.
\end{equation}
The program halts on success and halts after a failure with probability $\rho_*$. Its mean equals $N$, and the prediction risk is exactly $F/2-F^2/4$. The same program and forecast solve the robust problem with unknown fixed $h'\in[0,h]$.

For $0<a\le1/4$, $\sqrt a/2\le F\le\sqrt a$. Consequently the prediction-risk difference from $h=0$ is of order
$\sqrt{(N-1)\eta}$ with numerical constants, uniformly for $N\ge2$ in this small-error range. This matches the order in Theorem~\ref{thm:normalized-transfer}; no strictly larger exponent of the one-call perturbation can hold uniformly on its full class at this state budget.
\end{theorem}
\begin{proof}
With two total states there is one live state and one terminal state. A program is therefore specified by its stopping probabilities on $s,f$ and its terminal forecast. Halting immediately on success preserves the terminal physical success event and decreases the mean. Since squared risk optimized over the one forecast is increasing in the failure probability, this reduction is valid for both tasks, and includes all competitors.

Write $\rho>0$ for the failure stopping probability and
$D=p+q\rho+h(1-\rho)$. Direct summation of the two live occupations gives
\begin{align}
 z_A&=D^{-1},&z_T&=\frac{(1-\rho)h}{\rho D},\notag\\
 T(\rho)&=\frac{1+h(1-\rho)/\rho}{D},&
 F(\rho)&=\frac{q\rho+h(1-\rho)}{D}.
 \label{eq:instrument-occupations}
\end{align}
The terminal masses at $S,A,T$ are $p/D$, $\rho(q-h)/D$, $h/D$, respectively, and sum to one. These identities use the actual unconditioned transition probabilities. As $\rho$ increases from zero to one, the mean decreases continuously from infinity to one and $F(\rho)$ increases strictly; the latter follows from $F'(\rho)=p(q-h)/D^2>0$. Mean monotonicity also follows by coupling the failure-stop coins, and is strict here. Hence the optimum has $T(\rho)=1/p$.

Using $1-F=p/D$, this equality is
$F=h(1-\rho)(1-F)/\rho$. Substitute it into $D$ to obtain
$F^2=a(1-F)$, proving both formulas in \eqref{eq:two-state-instrument}. The prediction identity follows as in Theorem~\ref{thm:rare-memory}. For a fixed $\rho$, both $T(\rho)$ and $F(\rho)$ are increasing in $h$; explicitly their derivatives in $h$ are $p(1-\rho)^2/(\rho D^2)$ and $p(1-\rho)/D^2$. The endpoint program is therefore feasible and least favourable throughout the unknown parameter interval.

The raw change in the failure instrument moves mass $h$ from $A$ to $T$, giving induced $\ell^1$ norm $2h$; other input columns and reports are unchanged. The inequality $F^2=a(1-F)$ gives $F\le\sqrt a$, while $a\le1/4$ implies $F\le1/2$ and hence $F\ge\sqrt{a/2}\ge\sqrt a/2$. The prediction-risk comparison with $F$ proves the last claim.
\end{proof}

\subsection{Finite descriptions without budget enlargement}
The exact frontier concerns Borel programs, and its optimal probabilities may be irrational. This is separate from finite description. For rational $p,b,N$ with $b>0$, the clock in Theorem~\ref{thm:rare-memory} has certified finite approximants on the same states. Indeed $m_*\in[1,N/b]$, and the increasing function in \eqref{eq:rare-root} is explicitly computable with rational values at rational arguments on each branch. Bisection gives an interval enclosing $m_*$ to width at most $v/2$. Choose its lower endpoint $m$. If a phase boundary $k$ lies in the interval, it may first be split there. If $m\le k$, round the last continuation probability downward; if $m\ge k$, round the phase-advance probability $k/m$ upward. Either operation decreases the all-failure clock in stochastic order, so it decreases the actual mean. No comparison of an undecidable numerical boundary is required.

To make the additional rounding change in $m$ at most $v/2$, a dyadic mesh $Q^{-1}$ is sufficient if
\[
 Q^{-1}\le (v/2)\min\{1,kb^2/N^2\}.
\]
On the second branch $\rho=k/m\ge kb/N$ and
$k/\rho-k/(\rho+Q^{-1})\le N^2/(kb^2Q)$; the first branch is immediate. Since $g_{k,p}$ is $p$-Lipschitz, a total reduction of $m$ by $v$ raises failure risk by at most $pv$. Approximating the chosen constant forecast to accuracy $v$ adds at most $v^2$ to its optimal squared risk. The table uses at most $M$ persistent labels, a stored probability and readout, and a finite list of phase successors. Program bits and temporary coin-generation workspace depend on $\log Q$ and the input description, and are not included for free in $M$. Sampling still uses at most mean $N$ calls; preparation and the terminal audit are additional physical operations. These statements give sufficient description and execution bounds, not a lower frontier for all computational resources.

\section{Same-task error composition and paid implementation}\label{sec:errors}
First-discrepancy coupling, state-product simulation, and attained hidden-offset and erased-bit constructions are inherited mechanisms~\cite{GTFStopped,GTFBayes}. We state their precise form for the present mean-budget architecture, and combine them with its correlated acquired-geometry law.
\subsection{Stopped-law perturbations}
A numerical or calibration comparison must include the stopping decision and the protected audit. Approximation of the report marginal alone does not compare the scored task.

\begin{proposition}[Acquired-occupation error composition]\label{prop:stopped-error}
Consider two causal implementations on the same visible interface. Suppose their next conditional interface kernels, on a common pre-acquisition history $h$, admit a coupling with mismatch probability at most $\kappa(h)$; the protected final audit coupling has mismatch at most $\kappa_f(h)$. Use the same observer and fresh randomization in both. If the ideal implementation has finite mean stopping time, then the total variation distance of their complete stopped, scored laws is at most
\begin{equation}\label{eq:stopped-tv}
 \min\left\{1,\ \E\sum_{t=1}^{\tau}\kappa(H_{t-1})
                    +\E\kappa_f(H_\tau)\right\}.
\end{equation}
A possible nonterminating approximate run is recorded at an extra failure outcome. If $0\le\ell\le L$, the corresponding risk difference is at most $L$ times \eqref{eq:stopped-tv}. In particular, uniform bounds $\kappa\le\theta$, $\kappa_f\le\theta_f$ and $\E\tau\le N$ give $L(N\theta+\theta_f)$, truncated at $L$.
\end{proposition}
\begin{proof}
Couple the conditional kernels sequentially while the histories agree, and use identical routing coins. The event that the first disagreement occurs at acquisition $t$ has probability at most the ideal expectation of $\one_{\{\tau\ge t\}}\kappa(H_{t-1})$. After the ideal stop, couple the protected audit in the same way. The union bound and monotone convergence give \eqref{eq:stopped-tv}; on its complement the complete finite stopped records and audit outcomes agree. Coupling bounds total variation with convention $\norm{P-Q}_\TV=\sup_A|P(A)-Q(A)|$. Integrating a loss in $[0,L]$ gives the risk bound. This argument concerns the conditional kernels of the full compared interface, including the hidden-state effect needed for later reports or its equivalent visible conditional law. A bound on a single current-report marginal is not substituted for that hypothesis.
\end{proof}

\begin{remark}[Mean cost is a separate obligation]\label{rem:mean-not-tv}
Total variation does not control expected acquisition. Replacing a one-call stop by a branch of probability $\eps$ lasting $\eps^{-2}$ calls changes the complete-law probability of that branch by at most $\eps$, while adding order $\eps^{-1}$ to the mean. Thus \eqref{eq:stopped-tv} is not a budget certificate for the approximate machine. For the rejection construction one may plan an ideal conditional stop floor $(1+\zeta)/N$ and require the implemented acceptance probability to differ by at most $\zeta/N$ at every surviving history. The implemented floor is then $1/N$, which proves its own mean bound and geometric tail. Appendix~\ref{app:computation} gives a geometric alternative with a deterministic acceptance margin.
\end{remark}

If a fixed acceptance rule is retained but each final code vector is perturbed by at most $\delta_{\rm out}$, Minkowski's inequality gives
\begin{equation}\label{eq:root-error}
 \E\norm{X_\tau-\widehat r_J}^2
   \le\left(\sqrt{\E\norm{X_\tau-r_J}^2}+\delta_{\rm out}\right)^2.
\end{equation}
Only when $r_J=\E[X_\tau\mid J]$ under the same actual stopped law does orthogonality improve this to addition of $\delta_{\rm out}^2$. Centers optimized for a reference law need not be calibrated centroids for every unknown correlated model. The stronger additive formula is therefore not asserted for them.

\subsection{Resource-certified causal simulation}
\begin{proposition}[Complete-task causal transfer]\label{prop:morphism}
A parameter-independent causal simulator from a source experiment to a target experiment has at most $R$ internal states and a causal action lift. It transforms reports, stopping, and the protected task readout compatibly. Suppose that, uniformly over models and all target policies of interest, its complete stopped scored law has total variation error at most $\eps$, and its source acquisition certificate is
\[
 \E[\text{source calls}]\le\kappa\E[\text{target calls}]+c.
\]
Then a target observer with $M$ states and mean acquisition at most $N$ has a source implementation with at most $RM$ states and mean acquisition at most $\kappa N+c$, satisfying
\begin{equation}\label{eq:morphism-risk}
 \calR_{\rm source}(\kappa N+c,RM)
       \le\calR_{\rm target}(N,M)+L\eps.
\end{equation}
Two such simulators compose with error at most $\eps_1+\eps_2$, state multiplier $R_1R_2$, and acquisition bound
\begin{equation}\label{eq:morphism-cost}
 N\longmapsto\kappa_1\kappa_2N+\kappa_1c_2+c_1.
\end{equation}
Program descriptions, temporary workspaces and physical execution times remain separately charged.
\end{proposition}
\begin{proof}
Run the simulator and target observer in their product state. Only past reports and fresh coins enter the product transition. The action lift ensures that every requested target action has a legal source realization; compatibility of the stop and audit prevents information from arriving after the target decision but before the source decision. Apply the declared mean-cost certificate to the executed target policy and the total variation loss bound to the complete task. Infimizing over target observers proves \eqref{eq:morphism-risk}. Sequential composition gives \eqref{eq:morphism-cost} by substitution. Kernel contraction of total variation and the triangle inequality give the additive error. The product of internal state sets gives the state count. Reverse lower bounds require a reverse simulator or an independent lower argument; they do not follow from this one-sided construction.
\end{proof}

\subsection{An attained joint resource--risk law}
It is useful to distinguish an allowable error upper bound from an experiment that actually has that error. We give the latter in the same task. Start with a regular-support report-recoverable experiment from Theorem~\ref{thm:main}. A stopped core state is frozen. The protocol then requires exactly one bit acquisition: for an unknown model bit $b\in\{-1,1\}$ it reports $b$ with probability $1-e$ and an erasure symbol with probability $e$. The bit affects no preceding core report. An unknown calibration parameter $c\in\{-\delta,\delta\}$ affects no observer report at all. The final protected audit returns the core audit, $b$, and $c$, and the loss is the sum of their squared prediction losses. The observer must commit all three predictions before this audit. It is forbidden to restart the core or repeat the bit acquisition.

\begin{lemma}[Attained erasure deficiency]\label{lem:erasure}
For this bit interface, relative to the perfect bit experiment, the minimum parameter-independent causal simulation error is exactly $\eps=e/2$. The protected truth readout is fixed and cannot be altered by the simulator. The minimax bit prediction risk is $e$, and the independent hidden-calibration prediction risk is $\delta^2$.
\end{lemma}
\begin{proof}
On a nonerased report, output the reported bit. On erasure, a fair guess gives error $e/2$ for each model bit, so it achieves the stated total variation defect against the perfect report and unchanged protected truth. Conversely, under the uniform prior on $b=\pm1$, no pre-audit information on an erased branch depends on $b$. Every simulator then guesses wrongly with probability at least $1/2$ on that branch. Its largest model error is at least its prior average, $e/2$. This proves exact deficiency, without access to the final audit. For squared bit prediction, the conditional Bayes estimator on erasure is zero and its loss is one; the minimax risk is $e$, attained by zero on erasure and the reported sign otherwise. Since the calibration sign is absent from every report, averaging its two signs gives $\E(r-c)^2=\E r^2+\delta^2$, minimized at zero.
\end{proof}

\begin{theorem}[Attained acquisition--memory--error composition]\label{thm:joint-law}
Let $0\le\delta\le1$, $0\le\eps\le1/2$, $n\ge1$ and $M\ge5$. Give the core a mean acquisition budget $n$ and use the extended protocol with erasure probability $2\eps$. Its model-uniform optimum satisfies
\begin{equation}\label{eq:joint-law}
 \calR_{\rm ext}(n,M,\delta,\eps)-B
        \asymp (nM)^{-2/s}+\delta^2+\eps.
\end{equation}
The upper is implemented by at most $4m+1\le M$ states, with $m=\lfloor(M-1)/4\rfloor$. Under the convention charging one preparation, each core call, one bit call and one protected audit call, the mean total physical-call budget is $\mathsf N=n+3$. Thus the geometric term in that convention is $((\mathsf N-3)M)^{-2/s}$, not $(\mathsf NM)^{-2/s}$ with the extra calls silently omitted.
\end{theorem}
\begin{proof}
For the upper, use the one-live-state core construction with $m$ centers. Upon acceptance, retain its center index in one of $m$ bit-waiting states. These states perform the single bit call. Its three possible reports lead to $3m$ terminal states, carrying the fixed core center and the bit prediction in $\{-1,0,1\}$. The calibration prediction is zero. The total is $1+m+3m=4m+1$. No clock, free retained bit or exact posterior is present. The core risk bound and Lemma~\ref{lem:erasure} give
\[
 \calR_{\rm ext}-B\le C(nm)^{-2/s}+\delta^2+2\eps.
\]
Since $m\ge M/8$ for $M\ge5$, this is the upper in \eqref{eq:joint-law}.

For the converse, any legal extended observer commits to a core stop before receiving the bit. Its actual stopped core law is bounded by $b_*n\nu$. The first component of its final decoder has at most $M$ possible values, so its core squared distortion is at least $\calQ_M^\nu(b_*n)\ge c(nM)^{-2/s}$. This remains true even though the later bit report can select among those values. Average over the two bit signs and two calibration signs, keeping any one core model fixed. These signs do not change the preceding core law. Conditional on an erased bit there is still no information about its sign; there is never pre-audit information about the calibration sign. The three squared losses add, so their averaged lower bound is $c(nM)^{-2/s}+2\eps+\delta^2$. The worst model risk is at least this average. The physical-call count is a direct count of the specified protocol. There is no conversion from total variation to an expected-cost bound in this proof.
\end{proof}

Theorem~\ref{thm:joint-law} supplies attained lower and upper error terms for a fixed enlarged task. Proposition~\ref{prop:morphism} is the general transfer statement. They are logically different results: an arbitrary permitted calibration tolerance or simulator error need not be attained by a given experiment and hence need not cause that much risk.

\subsection*{Same-budget primitive perturbations}
For common ordered carriers, Lemma~\ref{lem:resolvent-error} strengthens the bounded-task comparison by a separate occupation-weighted mean-cost estimate after program-side stopping regularization. Proposition~\ref{prop:same-budget} then transfers the optimum at the same $N,M$ under uniform raw instrument and preparation defects. It does not assert that total variation of an arbitrary stopped transcript controls its expected length. A simulator that changes the carrier, consumes extra native calls, or retains additional state must still satisfy Proposition~\ref{prop:morphism}; the stopping argument does not erase its cost.

\appendix
\section{Executable tests and measurable predictive quotients}\label{app:quotient}
We give the elementary finite-closure construction underlying the representation in Section~\ref{sec:experiments}. Fix one prepared model. Let $\mathcal V$ be the linear space generated by bounded executable future-test functions on the physical state, including the constant function, the audit moments, and all legal one-step report pullbacks. A test is a specified causal continuation with a bounded final report, not an arbitrarily attached coordinate. Suppose $\mathcal V$ is finite-dimensional and admits the positive instrument representation \eqref{eq:raw-instrument}. Choose a basis $1,f_1,\ldots,f_d$ and evaluate it on prepared posterior measures.

\begin{proposition}[Realization and minimality]\label{prop:quotient}
Two prepared posterior measures give the same value to every test in $\mathcal V$ if and only if their basis evaluation vectors agree. Finite future report-and-action laws and the protected audit descend to this evaluation. On positive evidence, its update is \eqref{eq:update}. The history-to-evaluation map is Borel and gives a standard Borel realization in the compact base $D$. Any Borel statistic of a standard Borel history that determines all these future tests factors this evaluation through a Borel map.

This does not assert that the image of every Borel history map is a Borel subset of $D$. When only an analytic image is established, $D$ is the declared Borel realization and probability laws are pushforwards into $D$; unused points are not declared acquired. In the two raw realizations the post-report evaluation images are compact and therefore Borel.
\end{proposition}
\begin{proof}
Equality on a finite basis is equivalent to equality on its span. Pullback closure means that the probability of a report event and the unnormalized numerator of every subsequent test are linear evaluations of the incoming vector. Dividing by its positive report likelihood gives \eqref{eq:update}; zero likelihood has zero positive numerator by strict positivity of $e$ on the cone. Iterating gives every finite continuation law. Bounded monotone limits of such tests remain affine when they belong to the declared closure: evaluation on a fixed finite spanning family of prepared states determines the coefficients, so pointwise bounded convergence determines their limit. Alternatively, one may retain just the explicitly declared finite continuations; no closure under an unproved infinite-horizon operation is needed for the main theorem.

The posterior evaluation is Borel by standard Borel conditional kernels, or directly by finite matrix products and normalization on the finite-report history spaces, pieced together over their countable union. If another statistic $r(h)$ determines all tests, each evaluation coordinate is constant on its fibers. For completeness, the Borel factor statement follows by separating the two disjoint analytic images under $r$ of the Borel sets where a coordinate is, respectively, at most or greater than a fixed rational number. This shows that each coordinate is measurable with respect to $\sigma(r)$. The measurable factorization lemma for a real-valued function, applied coordinatewise, supplies a Borel factor; the resulting factor is defined on the entire standard Borel codomain of $r$. Its inverse image of the closed set $D$ is Borel and contains every attained statistic; replacing its values on the complement of that Borel set by a fixed point of $D$ makes it $D$-valued. No modification on a potentially non-Borel image complement is used. Thus the evaluation is the minimal information required by these tests, in this precise factorization sense.

The image caveat follows because Borel images of standard Borel spaces are in general only analytic. The general construction does not require the image to be the state space of the machine. In Theorem~\ref{thm:classical}, the map $x\mapsto\phi(x)$ is continuous and injective on compact $K$, and the audit probabilities separate the full convex base. In Theorem~\ref{thm:quantum}, the density matrices form a compact set and the informationally complete audit separates them. Their acquired rank orbits are also compact. These give the stronger image assertion for the realizations.
\end{proof}

For a family of unknown models, this construction is made within each model's finite predictive representation. A posterior on an unknown continuous model parameter can require additional, even infinite-dimensional, coordinates. Theorem~\ref{thm:main}(i) does not pretend otherwise: it couples the model-indexed certificates by a single observer kernel. It never gives that kernel access to a model-specific posterior. The geometric upper is stronger in a different respect: the current response $f_a(y)$ and the acceptance rule are common to the whole family.

\begin{example}[Separate feasibility does not give a common program]\label{ex:model-routing}
Prepare a fair bit $U$. In model $0$ the report is $Y=U$; in model $1$ it is $Y=1-U$. The protected audit is $U$. Both models have the same fair report law, and each separately has a perfect one-call decoder. No common decoder can be perfect: under the uniform prior on the two model labels, $U$ and $Y$ are independent fair bits. For squared $0/1$ prediction the common minimax risk is $1/4$, attained by reporting $1/2$. Maximizing the routing contribution separately within each model would incorrectly certify zero. This is why model contributions in \eqref{eq:joint-support} are added before its maximum, and why terminal readouts remain common.
\end{example}

\section{Finite descriptions and numerical resource certificates}\label{app:computation}
The exact optimization in Theorem~\ref{thm:main}(i) is over Borel kernels. It is a causal realization theorem in the stochastic-kernel model, not a theorem that every such kernel has a finite Turing description. Theorem~\ref{thm:effective-budget} supplies a general finite-description completion under effective raw data. Here we retain an independent, much more explicit finite-description upper for the sharp regular-support class. Its extra input is a certified numerical report and a computable geometric net, not a stored exact posterior.

\begin{proposition}[Finite-bit regular-support implementation]\label{prop:finite-bit}
Assume Definition~\ref{def:recoverable} and \eqref{eq:regularity}. Suppose the following data are supplied by the raw implementation: at accuracy $2^{-p}$, the current task response in $\R^k$ can be read with deterministic Euclidean error at most $C_y2^{-p}$ using workspace $W_y(p)$ and time $T_y(p)$; finite rational global covers and separated center lists at their asserted scales can be constructed; and the fixed action $a_0$ has a finite description of length $D_0$.

Then a finite program attains the upper order in \eqref{eq:regular-aut}, using one live label and $m=M-1$ terminal labels, with
\begin{equation}\label{eq:precision-cost}
 p\le C+\left\lceil\frac1s\log_2(Nm)\right\rceil,
 \qquad D_{\rm prog}\le D_0+Ckm(p+\log(m+1)).
\end{equation}
One valid temporary workspace bound is
\begin{equation}\label{eq:workspace-cost}
 W_{\rm temp}\le W_y(p)+Ck(p+\log(k+1))+C\log(m+1),
\end{equation}
and each live transition uses at most $T_y(p)+Cm\,T_{\rm dist}(k,p)$ arithmetic time. The constants include the fixed geometry and certified arithmetic implementation. The mean number of transitions is at most $N$; there is no deterministic upper bound on the number of transitions. The program and report computation do not access the unknown model.
\end{proposition}
\begin{proof}
For sufficiently large $N$ take $m$ separated centers as in Lemma~\ref{lem:regular}. Set $r=C_2(Nm)^{-1/s}$, where $C_2$ is a fixed large constant. Choose the large-$N$ threshold so that $3r$ is less than their separation. Round the centers and read the current response accurately enough that their combined distance error, including certified arithmetic, is at most $r/4$. Accept when the computed distance to some rounded center is at most $r$. The true union of radius-$r/2$ balls is then contained in the accepting set, and every accepted point is within $2r$ of its selected rounded decoder. The lower acquired mass bound implies, at every surviving history,
\[
 \Pp(\text{accept}\mid h)\ge a_*mc_0(r/2)^s\ge 1/N
\]
by the choice of $C_2$. Thus the \emph{implemented} stopping time has mean at most $N$, independently of any total variation approximation. Its squared distortion is at most $4r^2$.

For $N$ below the fixed threshold, use a certified global cover with $m$ centers, stop after the first report, and select a center of approximately smallest distance. The cover and numerical errors give $Cm^{-2/s}\le C'(Nm)^{-2/s}$ in this range. Both cases use $p=O(1+s^{-1}\log(Nm))$. Storing $m$ rounded $k$-vectors, a fixed threshold and their indices gives \eqref{eq:precision-cost}, allowing harmless additional index-address bits. Computing distances by scanning the read-only list requires one current response, one running best distance, one index, and arithmetic scratch space, which gives \eqref{eq:workspace-cost}. This workspace is erased before a rejected call returns to the single live label. The terminal center index is the only retained outcome. Time is counted for every attempted call, including failures.
\end{proof}

\subsection*{Certified nets in the two realizations}
For the box, rational lattice covers and separated sublattices provide the required lists. For the product Cantor set, depth-$h$ cylinders have rational endpoints, known mass $2^{-dh}$ and known separation at their natural scale. Taking enough endpoints supplies separated lists, and complete cylinder families supply covers. The relation $3^{-h}\asymp m^{-1/s}$, with $2^{dh}$ within a fixed factor of $m$, gives the asserted cardinality bounds. No membership oracle for an arbitrary null set is assumed.

For the quantum support, cover projective space by the $q$ charts in which $|\psi_j|$ is largest. After fixing that coordinate's phase, write the ray as $(1,z)/\norm{(1,z)}$, with all real and imaginary components of $z$ in $[-1,1]$. The corresponding projection
\[
 P(z)=\frac{(1,z)(1,z)^*}{1+\norm z^2}
\]
is Lipschitz on each compact chart with a constant depending only on $q$. Rational grids of mesh $h$ therefore give computable covers of cardinality $O_q(h^{-2(q-1)})$ and rational matrix coordinates. A greedy separated subset of a sufficiently fine such net supplies the packing lists: maximality gives a cover at a fixed multiple of the separation scale, and the exact mass formula \eqref{eq:projective-mass} bounds its cardinality above and below. Choosing the scale with a fixed safety factor produces at least $m$ separated points, or at most $m$ covering points, as needed. These are finite, verifiable constructions, not a claim of a fast optimal projective quantizer.

The literal description length in \eqref{eq:precision-cost} grows with precision even though the observer has only $M$ persistent labels. For the quantum instrument, its fixed physical dimension, channel invocation, outcome digitization and audit measurement are additional native costs. For the classical instrument, the vector-report acquisition and the scalar operations implementing it are likewise distinct. A bounded-state theorem alone cannot identify either of these total bit or physical cost frontiers.

\subsection*{What is not supplied by a finite check}
The support criterion \eqref{eq:joint-support} quantifies over real dual vectors and, for an arbitrary model family, all finite model subsets. The proof establishes the continuum statement by compactness and separation. Exact finite examples check special instances, signs and constants; they do not replace that proof. Nor does the geometric program calculate the optimal centers in \eqref{eq:profile}: it implements the proven matching order under explicit numerical certificates. Optimal program bits, workspace, planning time and simulator-state lower bounds remain separate mathematical questions.

\section{Description, workspace, and evaluation costs of the general construction}\label{app:finite-description}
This appendix accounts for the finite program in Theorem~\ref{thm:effective-budget}. It supplements, rather than replaces, the sharper geometric implementation in Appendix~\ref{app:computation}.

Let $l=|S|$, $j=|J|$, $a=|\calA|$, $m=l+j\le M$, and let every action partition have at most $b$ cells. Let grid probabilities have denominator $Q=2^q$, and let $\rho=2^{-k}$. More general dyadic $\rho$ are handled by the same finite-bit sampling. Suppose each legal decoder has a supplied description of at most $D_H$ bits, and the action and partition primitives have total description $D_{\rm raw}$.

\begin{proposition}[A literal finite-state implementation]\label{prop:description}
One explicit representation of the returned program uses
\begin{equation}\label{eq:program-bits}
 D_{\rm prog}\le D_{\rm raw}
  +C\bigl[la(q+1)+labm(q+1)+jD_H+k+m\log(m+a+b+1)\bigr]
\end{equation}
read-only program bits and at most $M$ persistent labels. If partition evaluation uses temporary workspace $W_h$ and time $T_h$, table scanning and exact dyadic sampling give temporary workspace
\begin{equation}\label{eq:general-workspace}
 W_{\rm temp}\le W_h+C\bigl(q+k+\log(m+a+b+1)\bigr)
\end{equation}
and per-acquisition transition time at most
\begin{equation}\label{eq:general-time}
 T_h+C(a+m)(q+k+\log(m+a+b+1))
\end{equation}
in a bit-operation model with sequential read-only table access to the selected row. More expensive table addressing may instead be charged explicitly; it is not a persistent-state resource. The expected number of physical acquisitions is at most $N$, and the protected audit is an additional specified physical call. There is no pathwise acquisition cap.
\end{proposition}
\begin{proof}
Store one dyadic action simplex per live label, one dyadic destination simplex per live-label/action/cell triple, all terminal decoder descriptions, the stopping parameter, and the finite row addresses. This gives \eqref{eq:program-bits}. To sample a dyadic row draw exactly $q$ fresh bits, interpret them as an integer in $\{0,\ldots,Q-1\}$, and scan the cumulative integer weights. If the destination is live, draw $k$ further fresh bits and replace it by the existing terminal label $j_0$ precisely when they are all zero. These coins are not retained. The report cell and all arithmetic are temporary and are erased at the transition boundary. Integer accumulation and row scanning give the stated workspace and time. No posterior vector, model index, persistent random seed, or horizon counter is present. The certified mean bound is \eqref{eq:exact-budget-certificate}; counting the protected audit separately does not alter it.
\end{proof}

For a decoder list of size $H_0$, a crude bound on the number of enumerated tables for fixed cardinalities is
\begin{equation}\label{eq:search-count}
 (Q+1)^{la+labm}H_0^j.
\end{equation}
Sum this over the finitely many cardinality choices. For each table and each model-net point, the transient ensemble dimension is $l\dim E$. Truncating after $t$ powers with
\[
 (1-\rho)^t/\rho\le\eta_T,\qquad L(1-\rho)^t\le\eta_R
\]
provides explicit tail intervals for mean and risk. Matrix products are evaluated to additional certified precision. These operations may be extremely costly as $\rho\downarrow0$. They are counted as offline synthesis and are not charged to the acquisition counter or silently removed from the total computational resource vector.

The bounds above are constructive upper bounds. Neither the program-description rate nor the exponential search count is claimed optimal. In particular, the existence of finite certificates does not supply an information-theoretic lower bound for all programs, workspace, and runtime simultaneously. The same-budget risk lower in Theorem~\ref{thm:effective-budget} concerns risk against all legal Borel observers, and is logically different from a lower bound on computational effort.

\section{Connection with informative Bayesian compression}\label{app:bayesian}
The sharp free-stopping profile is not a replacement for the more general informative-report identities. We record the latter to specify exactly what is inherited and what extra hypothesis the profile uses. The identity below is the normalized-pooling result of the separately retained Bayesian-curvature article~\cite{GTFBayes}; its short proof is included for independence of notation.

Let $U(y,x)$ and $\ell(y,x)$ be affine in a finite predictive vector $x$, with $\ell\ge0$, normalized report law and bounded posterior $F=U/\ell$ on positive evidence. Let $X$ be the exact incoming vector under an actually prepared law, $S$ its retained label, and $V=\E[X\mid S]$. Write $X'$ for the exact next vector, $Y$ for the current report and $J$ for any randomized outgoing label using only $(S,Y)$ and fresh coins. Define
\[
 \Psi(x)=\int\frac{\norm{U(y,x)}^2}{\ell(y,x)}\lambda(dy),\quad
 W=F(Y,V),\quad c_{sj}=\E[W\mid S=s,J=j],\quad v_j=\E[W\mid J=j],
\]
with zero contribution from zero-evidence reports.

\begin{proposition}[Inherited three-defect identity]\label{prop:bayesian}
Under the displayed raw affine and measurability assumptions,
\begin{align}
 \E\norm{X'-\E[X'\mid J]}^2
 &=\E\Psi(X)-\E\Psi(V)\notag\\
 &\quad+\E\norm{W-c_{SJ}}^2+\E\norm{c_{SJ}-v_J}^2.\label{eq:three-defects}
\end{align}
All three right-hand terms are nonnegative. The first equals
$\E\norm{X'-\E[X'\mid S,Y]}^2$.
\end{proposition}
\begin{proof}
Affine unnormalized integration gives the report density $\ell(y,V)$ conditional on $S$ and the corresponding numerator $U(y,V)$. Thus $\E[X'\mid S,Y]=W$, with the necessary likelihood reweighting already included. Fresh routing randomization gives $\E[X'\mid S,Y,J]=W$. Consequently $\E[X'\mid J]=v_J$. Two nested conditional squared-norm decompositions yield
\[
 \E\norm W^2=\E\norm{W-c_{SJ}}^2+\E\norm{c_{SJ}-v_J}^2+\E\norm{v_J}^2.
\]
Integrating the raw report kernels gives $\E\norm{X'}^2=\E\Psi(X)$ and $\E\norm W^2=\E\Psi(V)$. Subtract the last term of the preceding display to obtain \eqref{eq:three-defects}. Conditional projection proves the claimed form and sign of the first difference. No differentiability is needed.
\end{proof}

Rowwise optimal quantization does not imply equality of the pooled centroids $c_{sj}$, and one-step equality does not imply that compatible rows can be chosen throughout an arbitrary future experiment. For a general filter, the occupation criterion of Theorem~\ref{thm:main}(i) enforces that a single whole-run observer exists. Recoverability is the additional raw condition that makes the current task vector $W=f_a(Y)$ known from the current report and permits the sharp free-stopping construction. The retained informative Gaussian and noncommuting weak-effect filtering proofs use a different mechanism, namely actual mass, global covers and summable continuation sensitivity. They remain independent positive results; their fixed-task or prescribed-stop conclusions are not automatically statements about optimization under a free mean budget.

\section{A finite-state clock inequality}\label{app:clock-order}
This appendix records a standard phase-type comparison in the precise discrete form needed above. O'Cinneide~\cite{OC} proves the stronger convex-order minimality of an Erlang law of fixed phase order and mean, extending the coefficient-of-variation result of Aldous and Shepp~\cite{AS}. The following direct Laplace-transform proof is included to expose the state-count and time conventions; the phase-type extremality itself is not a new claim.

\begin{lemma}[Laplace bound for a finite absorption mechanism]\label{lem:phase-laplace}
Let $W$ be the almost surely finite absorption time of a continuous-time Markov chain with at most $k\ge1$ transient states, allowing an initial mixture and an atom at zero. Suppose $\E W=m<\infty$. Then for every $s\ge0$,
\begin{equation}\label{eq:phase-laplace}
 \E e^{-sW}\ge(1+sm/k)^{-k}.
\end{equation}
\end{lemma}
\begin{proof}
We prove the assertion by induction, first for a fixed initial state $i$. In a transient state the total occupation time, obtained by suppressing excursions away from that state, is exponential: its visits have independent exponential holding times, and the number of returning excursions is geometric. This proves the assertion for one transient state. Mixtures with an atom at zero are handled by convexity of $m\mapsto(1+sm)^{-1}$.

For the induction step, split $W=U+V$, where $U$ is time spent at the initial state and $V$ is time spent at the other states. Let $r<1$ be the probability of returning to $i$ after an exit. If $r=0$ the following formula is read with a geometric count identically zero. The strong Markov property gives a geometric number $K$ of returning excursions, with
\[
 \E z^K=g(z)=\frac{1-r}{1-rz},\qquad 0\le z\le1.
\]
Holding times at $i$ are exponential and independent of the destination and of the excursion. Conditional on a return, the excursion duration has one fixed law; conditional on the final nonreturn, its remaining duration has another fixed law. The successive pieces are independent conditional on $K$, with distributions not depending on its value. If $u$ is the holding-time Laplace transform, $v$ the final nonreturn duration transform and $y$ the returning-excursion transform, then
\[
 \E e^{-sW}=uvg(uy),\quad
 \E e^{-sU}=u g(u),\quad
 \E e^{-sV}=v g(y).
\]
For $0\le u,y\le1$ one has $g(uy)\ge g(u)g(y)$, because
\[
 (1-ru)(1-ry)-(1-r)(1-ruy)=r(1-u)(1-y)\ge0.
\]
Thus the Laplace transform of $W$ is at least the product of those of $U$ and $V$. The trace process obtained by suppressing time at $i$ is a continuous-time Markov chain on at most $k-1$ transient states, with a possible initial atom at absorption. This follows directly by summing the geometrically many instantaneous returns through $i$; equivalently its generator on the remaining states is the corresponding Schur complement. Its absorption time is $V$. The trace on $\{i\}$ has the exponential absorption time $U$. Put $a=\E U$ and $b=\E V$. By induction and the arithmetic--geometric mean inequality,
\[
 \E e^{-sW}\ge\frac1{1+sa}\frac1{(1+sb/(k-1))^{k-1}}
       \ge\frac1{(1+s(a+b)/k)^k}.
\]
Degenerate trace components follow by a limit or by deleting zero occupation states. Finally, for a mixture of initial states, apply the fixed-state conclusion and the convexity of $m\mapsto(1+sm/k)^{-k}$. This also includes the possible atom at zero.
\end{proof}

\begin{lemma}[Exact discrete clock frontier]\label{lem:discrete-clock}
Fix $p\in(0,1)$, $q=1-p$, and $k\ge1$. Let $T\ge1$ be the absorption-call count of any homogeneous randomized machine with at most $k$ live states and one terminal state. Assume $\E T=m<\infty$. Then
\begin{equation}\label{eq:clock-frontier-bound}
 \E q^T\ge g_{k,p}(m),
\end{equation}
where
\begin{equation}\label{eq:clock-frontier}
 g_{k,p}(m)=
 \begin{cases}
 (1-\beta)q^j+\beta q^{j+1},
    &1\le m\le k,\quad j=\lfloor m\rfloor,
      \ \beta=m-j,\\[2pt]
 (1+pm/(qk))^{-k},&m\ge k.
 \end{cases}
\end{equation}
The two expressions agree at $m=k$. Every point of this frontier is attained with at most $k$ live states. The function $g_{k,p}$ is continuous and strictly decreasing from $q$ to zero, and is Lipschitz with constant at most $p$ on $[1,\infty)$.
\end{lemma}
\begin{proof}
For $m\le k$, the piecewise linear interpolant of $j\mapsto q^j$ is convex. Jensen's inequality gives the first bound, attained by $T\in\{j,j+1\}$ with probabilities $1-\beta,\beta$. Such a clock uses $\lceil m\rceil$ live states: it advances deterministically until the $j$th call, halts there with probability $1-\beta$, and otherwise halts at the next call. If $\beta=0$ it uses only $j$ states.

For $m\ge k$, attach independent mean-one exponential holding times to the calls of the discrete machine, only for purposes of proof. The resulting time $W$ is phase type with at most $k$ transient states, generator $Q-I$ if $Q$ is the live routing matrix, and $\E W=\E T=m$. At $s=p/q$, conditioning on $T$ gives
$\E e^{-sW}=\E(1+s)^{-T}=\E q^T$. Lemma~\ref{lem:phase-laplace} proves the second bound. It is attained by the sum of $k$ independent geometric variables with success probability $\rho=k/m\le1$: a machine successively traverses $k$ phases, advancing one phase with probability $\rho$ at each call and halting on the advance from the last phase. Its mean is $k/\rho$ and its transform is
\[
 \left(\frac{q\rho}{p+q\rho}\right)^k
 =\left(1+\frac{pm}{qk}\right)^{-k}.
\]
This construction is discrete and pays one actual call for each step. The auxiliary exponential holds supply no physical clock or computational resource. On the linear pieces the derivative magnitude is $pq^j\le p$; on $m\ge k$ it is at most its value $pq^k\le p$ at $k$. These formulas prove all remaining assertions.
\end{proof}

\begin{thebibliography}{99}
\bibitem{CGM} J.~A.~Cuesta-Albertos, A.~Gordaliza and C.~Matr\'an, Trimmed $k$-means: an attempt to robustify quantizers, \emph{Ann. Statist.} \textbf{25} (1997), 553--576. \href{https://doi.org/10.1214/aos/1031833664}{doi:10.1214/aos/1031833664}.
\bibitem{DPR} F.~Dufour and T.~Prieto-Rumeau, Absorbing Markov decision processes, \emph{ESAIM Control Optim. Calc. Var.} \textbf{30} (2024), article 5. \href{https://doi.org/10.1051/cocv/2024002}{doi:10.1051/cocv/2024002}; also \href{https://arxiv.org/abs/2309.07059}{arXiv:2309.07059}.
\bibitem{FGPR} T.~Fritz, T.~Gonda, P.~Perrone and E.~F.~Rischel, Representable Markov categories and comparison of statistical experiments in categorical probability, \emph{Theoret. Comput. Sci.} \textbf{961} (2023), 113896. \href{https://doi.org/10.1016/j.tcs.2023.113896}{doi:10.1016/j.tcs.2023.113896}.
\bibitem{Kim} M.~Kim, Finite memory belief approximation for optimal control in partially observable Markov decision processes, preprint (2026), \href{https://arxiv.org/abs/2601.03132v1}{arXiv:2601.03132v1}.
\bibitem{PB} P.~Poupart and C.~Boutilier, Value-directed compression of POMDPs, \emph{Advances in Neural Information Processing Systems} \textbf{15} (2002), 1547--1554.
\bibitem{GTFBayes} Q.~Qi, General Theta Foundations I: acquired geometry and causal resource transfer, Bayesian-curvature technical article (2026). Complete retained mathematical source and native companion accompany this article.
\bibitem{GTFMean} Q.~Qi, General Theta Foundations I: acquired geometry and causal resource transfer, causal-stopping technical article (2026). The occupation realization, mean-constrained profile, and their complete proofs are retained with this article.
\bibitem{DPRG} F.~Dufour and T.~Prieto-Rumeau, Absorbing Markov decision processes: geometric properties and sufficiency of finite mixtures of deterministic policies, preprint (2025), \href{https://arxiv.org/abs/2512.17511}{arXiv:2512.17511}.
\bibitem{GTFStopped} Q.~Qi, General Theta Foundations I: acquired geometry and causal resource transfer, adapted-stopping technical article (2026), especially the independent-preparation frontier and the costed informative-instrument theorem. Complete retained mathematical source and native companion accompany this article.
\bibitem{SC} C.~R.~Shalizi and J.~P.~Crutchfield, Computational mechanics: pattern and prediction, structure and simplicity, \emph{J. Statist. Phys.} \textbf{104} (2001), 817--879. \href{https://doi.org/10.1023/A:1010388907793}{doi:10.1023/A:1010388907793}.
\bibitem{Sokolov} K.~Sokolov, Markovization of randomized stopping times, preprint (2026), \href{https://arxiv.org/abs/2609.24652v1}{arXiv:2609.24652v1}.
\bibitem{Varshney} K.~R.~Varshney, Bayes risk error is a Bregman divergence, \emph{IEEE Trans. Signal Process.} \textbf{59} (2011), 4470--4472. \href{https://doi.org/10.1109/TSP.2011.2159500}{doi:10.1109/TSP.2011.2159500}.
\bibitem{VPNC} A.~V\'erine, M.~S.~Pydi, B.~Negrevergne and Y.~Chevaleyre, Optimal budgeted rejection sampling for generative models, \emph{Proc. 27th International Conference on Artificial Intelligence and Statistics}, PMLR \textbf{238} (2024), 3367--3375.
\bibitem{GTFExactBudget} Q.~Qi, General Theta Foundations I: acquired geometry and causal resource transfer, exact-budget stopped-risk technical article (2026). Its complete mathematical source and companion accompany this article.
\bibitem{PZ} A.~Piunovskiy and Y.~Zhang, On the continuity of the projection mapping from strategic measures to occupation measures in absorbing Markov decision processes, preprint (2023), \href{https://arxiv.org/abs/2311.14043v1}{arXiv:2311.14043v1}.
\bibitem{AS} D.~Aldous and L.~Shepp, The least variable phase type distribution is Erlang, \emph{Comm. Statist. Stochastic Models} \textbf{3} (1987), 467--473. \href{https://doi.org/10.1080/15326348708807067}{doi:10.1080/15326348708807067}.
\bibitem{OC} C.~A.~O'Cinneide, Phase-type distributions and majorization, \emph{Ann. Appl. Probab.} \textbf{1} (1991), 219--227. \href{https://doi.org/10.1214/aoap/1177005935}{doi:10.1214/aoap/1177005935}.
\end{thebibliography}

\end{document}
